\documentclass[11pt]{article}
\usepackage[a4paper,margin=24mm,headheight=15pt]{geometry}
\usepackage[T1]{fontenc}
\usepackage{lmodern,amsmath,amssymb,amsthm,mathtools,booktabs,array}
\usepackage{xcolor,microtype,enumitem,fancyhdr}
\usepackage[colorlinks=true,linkcolor=blue!40!black,citecolor=blue!40!black,urlcolor=blue!40!black]{hyperref}
\hypersetup{pdftitle={Pilz's Extended 1-2-3 Conjecture over every field},pdfsubject={A support bound over every field}}
\pagestyle{fancy}\fancyhf{}
\lhead{Pilz's Extended 1--2--3 Conjecture over every field}\cfoot{\thepage}
\setlength{\emergencystretch}{3em}
\widowpenalty=10000\clubpenalty=10000
\setlist{itemsep=2pt,topsep=4pt}
\newcommand{\F}{\mathbb F_2}\newcommand{\N}{\mathbb N}
\DeclareMathOperator{\wt}{wt}\DeclareMathOperator{\ord}{ord}
\DeclareMathOperator{\supp}{supp}\DeclareMathOperator{\rank}{rank}
\DeclareMathOperator{\Span}{span}\DeclareMathOperator{\dist}{\delta}
\numberwithin{equation}{section}
\newtheorem{theorem}{Theorem}[section]
\newtheorem{lemma}[theorem]{Lemma}
\newtheorem{proposition}[theorem]{Proposition}
\newtheorem{corollary}[theorem]{Corollary}
\theoremstyle{definition}\newtheorem{definition}[theorem]{Definition}
\theoremstyle{remark}\newtheorem{remark}[theorem]{Remark}
\setcounter{tocdepth}{1}
\title{Pilz's Extended 1--2--3 Conjecture over every field\\[3pt]
\large Sparse multiples of weighted integer intervals}
\author{}
\date{23 October 2026}
\begin{document}
\maketitle
\begin{abstract}
We prove Pilz's Extended 1--2--3 Conjecture for every $n$, in a sharp
weighted form valid over every field. For every field $K$, every $n\ge1$
and all $c_1,\ldots,c_n\in K^\times$, every nonzero polynomial or
Laurent-polynomial multiple of $\sum_{a=1}^n c_a\prod_p x_p^{v_p(a)}$
has at least $n$ terms; Pilz's conjecture is the case $K=\mathbb F_2$,
$c_a=1$. The proof builds on the framework of Holdridge and Pach and adds
support compression and descent in widths and exponent residues, which
control cancellation over arbitrary fields, including imperfect ones.
It yields a general finite-residue formula for polynomials completed by
fresh variables. Further consequences include exact criteria for module
coefficients and commuting monoid actions, a coefficient-rank bound,
classification of equality at prime indices, and a universal inverse
after fewer than $n$ erasures. Every theorem stated here is
machine-checked in an accompanying Lean~4 development.
\end{abstract}

\begingroup
\small
\makeatletter
\renewcommand{\l@section}{\@dottedtocline{1}{0em}{1.8em}}
\makeatother
\tableofcontents
\endgroup

\section{Weighted integer intervals}\label{sec:main}

For each prime $p$, let $x_p$ be an independent variable, and write
\[
 X(a)=\prod_p x_p^{v_p(a)}\qquad(a\ge1).
\]
Unique factorization gives $X(ab)=X(a)X(b)$. Put $[n]=\{1,\ldots,n\}$.
All polynomials use finitely many variables. Write $\wt Q$ for
the number of nonzero monomials of a polynomial $Q$, with $\wt0=0$,
and, for $Q\ne0$, put
\[
 d(Q)=\min_{H\ne0}\wt(QH).
\]

\begin{theorem}[Weighted interval theorem]\label{thm:field}
For every field $K$, every $n\ge1$, and every
$c_1,\ldots,c_n\in K^\times$, the polynomial
\[
 P_{n,c}=\sum_{a=1}^n c_aX(a)
\]
satisfies $d(P_{n,c})=n$. The same bound holds for Laurent
polynomial multiples.
\end{theorem}

Equivalently, every nonzero finitely supported function
$f:\mathbb N_{>0}\to K$ satisfies
\begin{equation}\label{eq:convolution}
 (T_cf)(m)=\sum_{\substack{1\le a\le n\\a\mid m}}c_af(m/a),
 \qquad |\supp(T_cf)|\ge n.
\end{equation}
Indeed, the displayed sum is the coefficient of $X(m)$ in
$P_{n,c}\sum_b f(b)X(b)$. A singleton input gives equality.

The original question is the special case $K=\mathbb F_2$, $c_a=1$.
Over $\mathbb F_2$, the coefficient of $X(m)$ in $P_{n,1}F_A$, where
$F_A=\sum_{a\in A}X(a)$, is the parity of the number of
representations $m=ab$ with $a\in A$ and $1\le b\le n$.

\begin{corollary}[Pilz's Extended 1--2--3 Conjecture]\label{cor:pilz}
For every finite nonempty set $A$ of positive integers and every $n\ge1$,
\[
 \left|\mathop{\triangle}_{a\in A}\{a,2a,\ldots,na\}\right|
 \ge n .
\]
Equivalently, at least $n$ integers have an odd number of ordered
representations $ab$ with $a\in A$ and $1\le b\le n$.
\end{corollary}

Following~\cite{HP}, we write $A*[n]$ for the set in
Corollary~\ref{cor:pilz}, the integers with an odd number of such
representations. In the Lean development the corollary is
\path{WeightedPilz.weighted_symmetric_difference}; see the end of this
section.

\paragraph{A small example.}
With $x=X(2)$ and $y=X(3)$, the unweighted interval polynomial for
$n=4$ is $1+x+x^2+y$. Over every field,
\[
 (1+x+x^2+y)(1-x)=1-x^3+y-xy.
\]
Cancellation removes two intermediate powers of $x$, yet four terms
remain, attaining the bound. The nonmonomial cofactor also shows why
equality at a composite index can differ from the prime-index rigidity
in Corollary~\ref{cor:prime-rigidity}.

\begin{remark}[Scope of the bound]\label{rem:scope}
The interval structure is essential. Over $\mathbb F_2$,
\[
 (1+x+x^2)(1+x)=1+x^3,
\]
so the analogous bound fails for the three-element set $\{1,2,4\}$; this
example is from~\cite[\S1]{HP}.
The nonzero-weight condition is also exact: for arbitrary
$c_1,\ldots,c_n\in K$, the inequality
$\wt(P_{n,c}H)\ge n$ holds for every $H\ne0$ if and only if
every $c_a$ is nonzero. Sufficiency is Theorem~\ref{thm:field};
necessity follows by taking $H=1$, since the monomials $X(a)$ are distinct.
\end{remark}

\paragraph{Historical context.}
The question originates in Pilz's study of the minimal distance of
polynomial near-ring codes~\cite{Pilz}. Pilz proved the bound for
$n\le6$, proved the general lower bound $\pi(n)+2$, and reported that a
computer search found no counterexample.
Huang, Ke and Pilz~\cite{HKP} proved the interval-input case $A=[k]$,
known as the 1--2--3 Conjecture, and stated the general problem as the
Extended 1--2--3 Conjecture. Pach and Szab\'o~\cite{PS} gave the
number-theoretic formulation of Pilz's coding question, independently
proved the case $A=[k]$, and improved the general lower bound to
order $n/(\log n)^{\lambda}$ for every $\lambda>\lambda_0\approx0.2223$.
Ke and Meyer~\cite{KM} proved the conjecture for $|A|\le3$ and showed
that equality holds for arbitrarily large sets $A$, so the extremal
configurations are far from unique. Nagy and Pach~\cite{NP} proved the
analogous statement for additive sumsets. Most recently, Holdridge and
Pach~\cite{HP} proved the conjecture for $n\le8$ and for all
$n\ge2(3^{80}-321)$, and proved a lower bound of order $kn/\log k$ when
$k=|A|$ is small compared with $n$. Building on a formula of Ke and
Meyer, they also studied the limiting density
$h(A)=\lim_n|A*[n]|/n$, proving $h(A)\ge8/7$ unless $A=\{a\}$ or
$A=\{a,2a\}$. Theorem~\ref{thm:field} covers every $n$ and extends the
binary statement to arbitrary nonzero coefficients over every field.

This paper owes its foundations to Holdridge and Pach~\cite{HP}. The
encoding of integers by monomials in prime variables, the special role of
the primes in $(n/2,n]$, slicing by prime valuations, the decomposition
into smooth and rough components, and splitting estimates of the kind
used in Section~\ref{sec:large-weighted} and Appendix~\ref{sec:arithmetic}
all appear in their work, and the overall strategy for large $n$ follows
the path they opened.

\paragraph{What is new.}
The new ingredient is the treatment of multipliers that involve the
primes in $(n/2,n]$. Support compression, Frobenius extraction with
descent in base widths, and residue capacities (Sections~\ref{sec:tools}
and~\ref{sec:residue}) show that for $n\ge67$ some minimum multiple has a
cofactor that uses only the base variables, or, outside characteristic
two, has the affine form $D+ZB$ with $D$ and $B$ in the base variables.
After this reduction, component counting in the spirit of~\cite{HP}
closes the argument. These tools control cancellation uniformly in the
weights and over every field, and the uniform argument applies from
$n\ge67$, compared with $n\ge2(3^{80}-321)$ in~\cite{HP}; the
initial and middle ranges are treated separately in the appendices.

\subsection*{Outline of the proof}

A prime in $(n/2,n]$ occurs just once in the interval. Rescaling
these variables gives
\[
 P_{n,c}=C+Z,\qquad Z=z_1+\cdots+z_r,\qquad
 r=\pi(n)-\pi(n/2),\quad \wt C=n-r,
\]
where $\pi(t)$ counts primes at most $t$. The variables $z_i$
are fresh: none occurs in the core $C$. Each fresh variable records a
separate copy of a base-only multiplier, so prime factors absent from a
multiplier label groups of output terms that cannot cancel one another.
For $n\ge67$ the proof establishes three statements.

\begin{enumerate}[label=\textup{(\roman*)},leftmargin=2.2em]
\item \emph{Reduction} (Sections~\ref{sec:tools}--\ref{sec:large-weighted}). Some minimum multiple of $C+Z$ has
cofactor $D+ZB$ with $D,B$ in the base variables, and $B=0$ in
characteristic two.
\item \emph{Weight formula} (equation~\eqref{eq:large-affine}). Outside
characteristic two,
\[
 \wt\bigl((C+Z)(D+ZB)\bigr)=\wt(CD)+r\wt(D+CB)+\binom{r+1}{2}\wt B.
\]
\item \emph{Base-only bound} (equation~\eqref{eq:large-base}). For every
nonzero $D$ in the base variables,
\begin{equation}\label{eq:intro-base}
 \wt((C+Z)D)=\wt(CD)+r\wt D\ge n.
\end{equation}
\end{enumerate}

For (i), a finite support shift (Lemma~\ref{lem:two-variable-shift} and
Theorem~\ref{thm:compression}) combines the fresh variables into $Z$;
the aggregate degree of a multiple is its degree in $Z$. In positive
characteristic $p$, group these degrees modulo $p$: a correction
divisible by $(C+Z)^p$ can be extracted from each exponent residue class
(Lemma~\ref{lem:frobenius-extraction}). For a hypothetical counterexample
whose cofactor has least total width in the base variables
(Lemma~\ref{lem:width-descent}), every nonzero correction class has at
least the target number of terms, and bounds on the residue classes of
$C^j$ exclude these corrections (Theorem~\ref{thm:residue-descent}).
This gives a general finite-residue formula
(Corollary~\ref{cor:finite-residue}). For $n\ge67$ the fresh-prime
estimate $r\ge\lfloor\sqrt n\rfloor$ makes it apply uniformly, and a row of
aggregate degree three or more would already cost more than $n$ terms.

Given (i) and (ii), endpoint bounds show that $B\ne0$ already costs at
least $n$ terms, leaving~\eqref{eq:intro-base}. For (iii), delete terms
until no further deletion preserves a counterexample. If the multiplier
has $u$ terms, a split at any prime $q$ in its support forces
$q<2(u-1)$ (Lemma~\ref{lem:directed-cut}). Prime factors absent from the
multiplier then separate the output into disjoint components, and
counting these components (Theorem~\ref{arith:thm:localized-expansion})
excludes such a $q$.

\paragraph{The binary case.}
A reader interested only in Pilz's original statement, $K=\mathbb F_2$
and $c_a=1$, can follow a shorter path. For $n\ge67$ it uses, besides the basic tools of
Section~\ref{sec:tools}, Lemma~\ref{lem:two-variable-shift},
Theorem~\ref{thm:compression},
Lemmas~\ref{lem:frobenius-extraction} and~\ref{lem:width-descent},
Theorem~\ref{thm:residue-descent} in its characteristic-two form
(\S\ref{sec:residue}), and Corollary~\ref{cor:finite-residue} with $p=2$, which gives a
base-only cofactor; then~\eqref{eq:large-base}, via
Lemma~\ref{lem:directed-cut} and Theorem~\ref{arith:thm:localized-expansion}.
The weight formula in (ii) and the endpoint bounds are not needed.
For $n\le66$ it uses Appendices~\ref{sec:initial}
and~\ref{sec:finite}, specialised to characteristic two.

Section~\ref{sec:modules} gives the coefficient extensions: commuting
monoid endomorphisms have the same sharp support bound precisely when
each reflects zero. It also proves the rank bound $n+\rho-1$, classifies
equality at prime indices, and gives recovery after fewer than $n$
erasures.

\subsection*{Range and dependency guide}

\begin{center}
\begin{tabular}{@{}>{\raggedright\arraybackslash}p{.17\linewidth}>{\raggedright\arraybackslash}p{.77\linewidth}@{}}
\toprule
Range & Main mechanism \\
\midrule
$1\le n\le24$ & Initial cores, endpoint and binomial bounds, and special
cases (Appendix~\ref{sec:initial}). \\
$25\le n\le66$ & Finite-residue capacities and exceptional cases, with
a separate characteristic-two treatment (Appendix~\ref{sec:finite}). \\
$n\ge67$ & Uniform compression and descent, affine reduction, and
component counting (Section~\ref{sec:large-weighted}). \\
\bottomrule
\end{tabular}
\end{center}
The initial and middle ranges are proved uniformly over all fields,
nonzero weights, and nonzero multipliers. Their tables certify integer
inequalities used in these universal arguments, rather than testing a
finite collection of polynomials. The large-range proof uses only the
specified initial-case bounds and the independently proved arithmetic
estimates of Appendix~\ref{sec:arithmetic}; it does not assume the earlier
unweighted interval theorem.

\begingroup\raggedright
\paragraph{Formalisation and availability.}
This paper, its Lean~4 sources, and reproducible verification
instructions are archived on Zenodo~\cite{Archive}. The development uses
Lean~4.24.0 and mathlib revision
\texttt{f897ebcf72cd16f89ab4577d0c826cd14afaafc7}, and depends only on
the standard axioms \texttt{propext}, \texttt{Classical.choice} and
\texttt{Quot.sound}. Its main declarations are
\path{WeightedPilz.weighted_interval},
\path{WeightedPilz.weighted_interval_laurent},
\path{WeightedPilz.weighted_convolution}, and
\path{WeightedPilz.weighted_symmetric_difference}. The last is stated for
the set of Corollary~\ref{cor:pilz}, which is proved equal
(\path{intervalSymmetricDifference_eq_oddProduct}) to the odd product
\path{Pilz.oddProduct} of the earlier formal proof of the binary case.
The file \path{WeightedPilz/PROOF_MAP.md} in the archive records the
correspondence between this paper and the formal development. The
formal development verifies the mathematical statements represented
here, with explicitly documented alternative proofs of some intermediate
results; it does not claim a line-by-line translation of the printed
proof.

The same archive contains an earlier proof of the binary case
$K=\mathbb F_2$, $c_a=1$, which uses a square-root descent specific to
characteristic two and has its own Lean~4 formalisation, and a short
companion note proving that $h(A)\ge4/3$ unless $A=\{a\}$ or
$A=\{a,2a\}$, the improvement anticipated by Holdridge and
Pach~\cite[Remark~1.9]{HP}; the constant $4/3$ is attained by
$A=\{1,3\}$~\cite[Lemma~5.6]{HP}. The weighted proof here does not use the earlier binary
theorem as an input, although the two formal developments share their
arithmetic estimates.
\par\endgroup

\section{Elementary support tools}
\label{sec:tools}

\subsection{Basic operations and widths}

Extending the coefficient field does not change $d(Q)$. For fixed
finite sets of allowed cofactor and output exponents, the required
cancellations form a
homogeneous linear system over the original field. It has a nonzero
solution over an extension exactly when it has one over the original
field, by equality of matrix ranks. Multiplication by $Q\ne0$ is
injective, so either solution gives a nonzero product.

A common monomial may be removed from a cofactor without changing
the weight of its product. Polynomial and Laurent distances therefore
agree. Variables absent from $Q$ may also be removed: a nonzero
coefficient slice of the cofactor gives a nonzero $Q$-multiple of no
greater weight.

For a finite nonempty exponent set $A\subseteq\mathbb Z^a$ and an
integer vector $v$, its width in direction $v$ is
$\max_{\alpha\in A}v\cdot\alpha-\min_{\alpha\in A}v\cdot\alpha$.
The width of a nonzero polynomial means the width of its support.
In particular $\operatorname{width}_{x_i}Q=
\deg_{x_i}Q-\operatorname{ord}_{x_i}Q$.
The highest and lowest occupied rows of a product are products of
the corresponding factor rows. They are nonzero because a polynomial
ring over a field is a domain. Consequently
\begin{equation}\label{eq:width-additivity}
 \operatorname{width}_v(QH)=
 \operatorname{width}_v(Q)+\operatorname{width}_v(H).
\end{equation}

\begin{lemma}[Endpoint and dimension bounds]\label{lem:endpoints}
Suppose $Q=\sum_{i=0}^aQ_iy^i$, where $a>0$ and $Q_0Q_a\ne0$.
Then
\[
 d(Q)\ge d(Q_0)+d(Q_a).
\]
If the support of $Q$ has affine dimension $s$ over $\mathbb Q$,
every nonzero multiple of $Q$ has at least $s+1$ terms.
In particular, a polynomial with a nonzero constant and a positive
pure power of each of $s$ variables has distance at least $s+1$.
\end{lemma}
\begin{proof}
The lowest and highest occupied $y$-rows of $QH$ are distinct,
and are nonzero multiples of $Q_0$ and $Q_a$ respectively. This
proves the first assertion.

For the second, let $L_Q$ and $L_{QH}$ be the rational spans of the
support differences. If an integer vector $v$ annihilates $L_{QH}$,
then $\operatorname{width}_v(QH)=0$. Equation~\eqref{eq:width-additivity}
gives $\operatorname{width}_v(Q)=0$, so $v$ also annihilates $L_Q$.
Ordinary rational linear algebra therefore gives $L_Q\subseteq L_{QH}$:
otherwise a rational separating linear functional, multiplied by a
common denominator, would contradict this implication.
Thus the support of $QH$ has affine dimension at least $s$ and
contains at least $s+1$ points.
\end{proof}

\begin{lemma}[Exponent dilation]\label{lem:exponent-dilation}
For every integer $q\ge1$,
\[
 d\bigl(Q(x_1^q,\ldots,x_a^q)\bigr)=d(Q).
\]
\end{lemma}
\begin{proof}
Decompose any multiplier uniquely as
\[
 H(x)=\sum_{0\le\alpha_i<q}x^\alpha H_\alpha(x_1^q,\ldots,x_a^q).
\]
Its product with $Q(x_1^q,\ldots,x_a^q)$ has disjoint exponent
residue classes, each nonzero class having the weight of the
$Q$-multiple $QH_\alpha$. This proves the lower bound. Substituting
$x_i^q$ into a multiple attaining $d(Q)$ proves the reverse bound.
\end{proof}

\begin{lemma}[Power gap]\label{lem:powergap}
Let $C$ be nonmonomial and let $\delta=d(C)$. In characteristic
zero, $\wt(C^j)\ge\delta+1$ for every $j\ge2$. In characteristic
$p>0$, the same conclusion holds whenever $j>0$ is not a power of $p$.
\end{lemma}
\begin{proof}
In characteristic zero, put $D=C$ and $h=j$. In characteristic $p$,
remove a common monomial and the largest common exponent power $p^s$
to write
\[
 C=x^\mu D(x_1^{p^s},\ldots,x_a^{p^s}),
\]
where two exponents of $D$ differ modulo $p$. Write $j=p^eh$ with
$p\nmid h$. Since $j$ is not a power of $p$, we have $h\ge2$.
Exponent dilation and Frobenius give $d(D)=\delta$ and
$\wt(D^h)=\wt(C^j)$.

In either case, let $\beta$ be the lexicographically largest exponent
of $D$. There is a coordinate $i$ such that
$(E_i-\beta_i)D\ne0$, where $E_i=x_i\partial/\partial x_i$.
Then
\[
 (E_i-h\beta_i)D^h=hD^{h-1}(E_i-\beta_i)D
\]
is a nonzero $D$-multiple. The operator creates no monomials and
deletes the leading term at exponent $h\beta$. Its weight is at
least $\delta$ and at most $\wt(D^h)-1$, as required.
\end{proof}

If $j>0$ and $r\ge1$, define
\[
 W_r(j)=\wt((z_1+\cdots+z_r)^j),\qquad W_r(0)=1.
\]
In characteristic zero,
\[
 W_r(j)=\binom{r+j-1}{j}.
\]
In characteristic $p$, write $j=\sum_e j_ep^e$, $0\le j_e<p$.
Frobenius and the unique digit expansion of each exponent give
\begin{equation}\label{eq:digitcost}
 W_r(j)=\prod_e\binom{r+j_e-1}{j_e}.
\end{equation}
In particular $W_r(j)\ge r$ for $j>0$. If $r\ge2$, equality occurs
only at $j=1$ in characteristic zero and at powers of $p$ in
characteristic $p$. Every other positive degree costs at least
$\binom{r+1}{2}$. Different aggregate degrees have disjoint expanded
supports.


\subsection{Adding fresh variables}
\begin{theorem}[Sharp completion, including rigidity]\label{thm:completion}
Let $0\ne C\in K[x]$, $c=\wt C$, $\delta=d(C)$, and let
$z_1,\ldots,z_r$ be independent fresh variables. Then
\[
 d(C+z_1+\cdots+z_r)\ge\min\{c+r,\delta+2r\}.
\]
If $r>c-\delta$ and $c+r\ge3$, every minimum multiple is a nonzero
scalar monomial times $(C+z_1+\cdots+z_r)^q$, where $q=1$ in
characteristic zero and $q$ is a characteristic power otherwise.
\end{theorem}
\begin{proof}
For one new variable, remove the lowest power of $z$ from a nonzero
multiple $G$; divisibility survives since $C+z$ and $z$ are coprime.
Its constant row $A_0\ne0$ is divisible by $C$, by substituting $z=-C$.
If $\wt G<\delta+2$, exactly one positive-degree term remains, say
$Mz^j$ with $M$ scalar monomial. Thus
\[
 G=M\bigl(z^j-(-C)^j\bigr).
\]
The power gap excludes non-Frobenius exponents in this weight range.
The remaining expressions have weight $c+1$. A monomial $C$ is the
elementary binomial case. Hence $d(C+z)\ge\min(c+1,\delta+2)$, and
iteration gives the bound.

At the strict threshold, the polynomial obtained before adjoining
the final fresh variable already has distance equal to its weight
$c+r-1\ge2$. A minimum multiple after adjoining that variable has
weight $(c+r-1)+1<(c+r-1)+2$. The one-variable argument gives its
stated form, including the removed monomial factor.
\end{proof}


\subsection{Combining the fresh variables}

Put $Z=z_1+\cdots+z_r$. We shall show that a sparsest multiple of
$C+Z$ can be chosen to depend on the fresh variables only through
$Z$. All weights still count monomials in the individual $z_i$.
The proof moves allowed monomials towards earlier fresh variables
until no further move is possible, then extracts a coefficient in
the coordinates $z_1,\ldots,z_{r-1},Z$.

\begin{lemma}[A two-variable support shift]\label{lem:two-variable-shift}
Fix $d\ge0$, write $m_k=X^kY^{d-k}$, and let $A\subseteq\{0,\ldots,d\}$.
Apply the substitution $T_t:Y\mapsto Y+tX$, with $X$ fixed.
There is a set $B$ of the same size as $A$ with the following properties.
\begin{enumerate}
\item If $a_1>\cdots>a_s$ and $b_1>\cdots>b_s$ list $A$ and $B$,
then $b_i\ge a_i$ for every $i$.
\item The highest coefficient in $t$ of any nonzero polynomial vector
in $T_t\langle m_a:a\in A\rangle$ belongs to
$\langle m_b:b\in B\rangle$.
\item If $B=A$, the span of the $m_a$ is preserved by $T_t$.
\end{enumerate}
Membership in the shifted span is understood over $K(t)$.
\end{lemma}
\begin{proof}
Order the coordinates $m_d,\ldots,m_0$. Let $B$ be the pivot set
obtained by row reduction of the vectors $T_1m_a$. Since
\[
 T_tm_a=\sum_{k\ge a}\binom{d-a}{k-a}t^{k-a}m_k,
\]
a reduced basis $m_b+\sum_{k<b}\lambda_{b,k}m_k$ at $t=1$
gives the basis
\[
 f_b(t)=m_b+\sum_{k<b}\lambda_{b,k}t^{k-b}m_k
\]
at general $t$. Every nonpivot coefficient has negative degree in $t$.
In an expansion $F(t)=\sum_bh_b(t)f_b(t)$, the coefficient $h_b$
is exactly the pivot coordinate of $F$, and is therefore a polynomial
when $F$ is. The highest coefficient of $F$ is consequently supported
on the pivots. This proves the second assertion.

For each threshold $k$, the original columns indexed by $a\ge k$
remain independent on the rows indexed by exponents at least $k$:
the submatrix on their original indices is triangular with diagonal
$1$. Thus there are at least as many pivots above each threshold
as there are original indices, proving the first assertion.

Finally suppose $B=A$. Descend through the elements of $A$.
The largest-index vector $T_1m_a$ can have no term above $a$,
for that would give a larger pivot. After subtracting already
available larger monomials in $A$, the same argument applies to
each subsequent vector. Hence $T_1$ preserves the coordinate span.
All coefficients $\binom{d-a}{k-a}$ with $a\in A$ and $k\notin A$
are therefore zero in $K$, so the displayed formula proves
preservation for formal $t$ as well.
\end{proof}

\begin{theorem}[Scalar support compression]\label{thm:compression}
Let $C\in K[x_1,\ldots,x_a]$, let $Z=z_1+\cdots+z_r$ with $r\ge1$,
and put $P=C+Z$. Every nonzero $P$-multiple $G$ has a nonzero replacement
$F\in K[x,Z]$ divisible by $P$, with $\wt F\le\wt G$.
No base monomial occurring in $F$ is new. In particular,
\[
 \operatorname{width}_{x_i}(F)\le
 \operatorname{width}_{x_i}(G)\qquad(1\le i\le a).
\]
\end{theorem}
\begin{proof}
The case $r=1$ is immediate. A \emph{profile} records a base exponent
and a total fresh degree. Fill out the finitely many profiles
occurring in $G$ with all possible fresh monomials. Among the
coordinate sets $S$ having the same number of monomials in each
profile as $G$, and whose span contains a nonzero $P$-multiple,
choose one maximizing
\[
 \Phi(S)=\sum_{x^\gamma z^\beta\in S}
                 \sum_{i=1}^r(r-i)\beta_i.
\]
This is a maximum over a nonempty finite collection. Write $V$ for
the coordinate span of $S$, and choose $0\ne G_0\in V\cap(P)$.

Fix $i<j$. Partition $S$ into blocks with the same base exponent,
the same fresh exponents outside $i,j$, and the same sum of the
$i$th and $j$th exponents. Apply Lemma~\ref{lem:two-variable-shift}
to the substitution $T_t:z_j\mapsto z_j+tz_i$ in every block.
The resulting coordinate set $S'$ has the same profile counts.
It is still admissible: if $D_t$ scales $z_i$ by $1-t$, then
\[
 T_tD_t(P)=P,
\]
and $T_tD_t(G_0)$ is a nonzero polynomial in $t$, is divisible by
$P$, and belongs to $T_t(V)$ over $K(t)$. Its highest coefficient
in $t$ is therefore a nonzero $P$-multiple in the span of $S'$.

The pivot comparison gives $\Phi(S')\ge\Phi(S)$, strictly if any
block changes, because moving one exponent from $z_j$ to $z_i$
increases its weight by $j-i$. Maximality forces every block to
remain unchanged. The final assertion of the lemma shows that $V$
is preserved by every upper substitution $z_j\mapsto z_j+tz_i$.
It is also preserved by diagonal scalings.

Now use the coordinates $u_i=z_i$ for $i<r$, together with $Z$.
The substitution
\[
 u_i\mapsto\lambda_i u_i,\qquad Z\mapsto Z
\]
is a composition of the preceding substitutions and scalings:
in the original coordinates it sends $z_i$ to $\lambda_i z_i$
for $i<r$, and sends $z_r$ to
$z_r+\sum_{i<r}(1-\lambda_i)z_i$. It preserves both $V$ and $(P)$.
Expand
\[
 G_0=\sum_\alpha u^\alpha F_\alpha(x,Z).
\]
Taking a nonzero coefficient of the independent parameters
$\lambda_i$ shows that $u^\alpha F_\alpha\in V\cap(P)$ for some
$F_\alpha\ne0$. Substitution $Z=-C$ and cancellation of $u^\alpha$
in the polynomial domain give $F_\alpha(x,-C)=0$. Monic division
then shows that $P$ divides $F_\alpha$. Removing the monomial
$u^\alpha$ preserves weight and changes no base exponent, so
$F=F_\alpha$ has all the required properties.
\end{proof}

\subsection{Frobenius extraction and descent}

In characteristic $p$, taking a $p$th power dilates every exponent by
$p$. The next lemma reverses this operation on a residue class without
requiring roots of arbitrary coefficients.

\begin{lemma}[Frobenius extraction]
\label{lem:frobenius-extraction}
Let $K$ have characteristic $p>0$, and let $q$ be a power of $p$.
Let $P,E$ be polynomials over $K$ in any finite set of variables.
Suppose $0\ne U=P^qE$ and every
exponent of $U$ is congruent to $\alpha$ modulo $q$, with
$0\le\alpha_i<q$. Then there are polynomials $J$ and $F=PJ\ne0$ over $K$
such that
\[
 \alpha+q\,\supp F\subseteq\supp U,
 \qquad \alpha+q\,\supp J\subseteq\supp E.
\]
In particular $\wt F\le\wt U$ and
$\operatorname{width}_{x_i}F\le\operatorname{width}_{x_i}U/q$.
\end{lemma}
\begin{proof}
Multiplication by $P^q$ preserves exponent residues and is injective,
so every exponent of $E$ has the same residue $\alpha$.
The subfield $K^q=\{a^q:a\in K\}$ acts on the finite-dimensional
space spanned by the coefficients of $E$. All coefficients of $U$
lie in that space. Choose a $K^q$-linear functional to $K^q$ that
is nonzero on one coefficient of $U$, and apply it coefficientwise.
Since the coefficients of $P^q$ lie in $K^q$, the resulting nonzero
polynomial is still $P^q$ times the image of $E$, and has support
contained in $\supp U$.
The image of $E$ is $x^\alpha J^q$ for some $J$ over $K$: its
coefficients lie in $K^q$ and its exponents have residue $\alpha$.
Thus the image of $U$ is $x^\alpha(PJ)^q$. Take $F=PJ$.
\end{proof}

\begin{lemma}[Descent in base widths]\label{lem:width-descent}
Let $K$ have characteristic $p>0$, let $0\ne C\in K[x]$, and put
\[
 c=\wt C,\qquad Z=z_1+\cdots+z_r,\qquad P=C+Z,\qquad r\ge1.
\]
Fix an integer $1\le N\le c+r$. If a nonzero $P$-multiple with fewer than $N$
terms exists, choose $G=PH$ among them minimizing
\[
 b(H)=\sum_i\operatorname{width}_{x_i}H.
\]
Then $b(H)>0$, and $G$ may be chosen in $K[x,Z]$.
For this choice, let $0\ne U\in(P^p)$ satisfy, for every base variable,
\begin{equation}\label{eq:descent-exponent-bounds}
 \operatorname{ord}_{x_i}U\ge\operatorname{ord}_{x_i}G,
 \qquad
 \deg_{x_i}U\le\deg_{x_i}G+(p-1)\deg_{x_i}C.
\end{equation}
Every nonzero full exponent-residue class of $U$ modulo $p$ then
has at least $N$ terms.
\end{lemma}
\begin{proof}
Compression does not increase the weight or any base width of $G$.
By width additivity it does not increase any base width of $H$
either, so it preserves the minimum defining $b(H)$.

If $b(H)=0$, remove its base monomial to make $H$ depend only on
the fresh variables. Let $c_0$ be the constant coefficient of $C$,
and put $\epsilon=1$ if $c_0\ne0$, and $0$ otherwise. The disjoint
base rows give
\[
 \wt(PH)=(c-\epsilon)\wt H+\wt((c_0+Z)H)
 \ge(c-\epsilon)+(r+\epsilon)=c+r\ge N,
\]
a contradiction. Here $c_0+Z$ has distance $r+\epsilon$ by its
weight and the dimension bound in Lemma~\ref{lem:endpoints}.

Write $U=P^pE$. Multiplication by $P^p$ preserves full exponent
residues, so a nonzero class $U_\alpha$ equals $P^pE_\alpha$.
Lemma~\ref{lem:frobenius-extraction} gives a nonzero multiple $PJ$
with $\wt(PJ)\le\wt U_\alpha$ and base widths at most $1/p$ times
those of $U_\alpha$. Set $d_i=\deg_{x_i}C$. The fresh summand
ensures $\operatorname{ord}_{x_i}P=0$ and
$\operatorname{width}_{x_i}P=d_i$. Thus
\[
 d_i+\operatorname{width}_{x_i}J
 \le\frac{\operatorname{width}_{x_i}G+(p-1)d_i}{p}
 =d_i+\frac{\operatorname{width}_{x_i}H}{p}.
\]
Consequently $b(J)\le b(H)/p<b(H)$. Minimality gives
$\wt(PJ)\ge N$, and hence $\wt U_\alpha\ge N$.
\end{proof}

\subsection{A cut at a support prime}

Normalize a multiplier $D$ by removing its common monomial. Its
\emph{support primes} are the variables on which its exponents vary.
A counterexample is \emph{deletion minimal} if none of its nonzero
proper subpolynomials remains a counterexample.

\begin{lemma}[Directed cut]\label{lem:directed-cut}
Let $u=\wt D$, and suppose every nonzero proper subpolynomial of $D$
has product with $P_{n,c}$ of weight at least $n$. For any support
prime $q$,
\[
 \wt(P_{n,c}D)\ge2n-2(u-1)\left\lfloor\frac nq\right\rfloor.
\]
In particular, a deletion-minimal counterexample satisfies
$q<2(u-1)$.
\end{lemma}
\begin{proof}
Write $D=L+H$, where $L$ has zero $q$-valuation and $H$ has
$q$-valuation at least $b>0$. Both products have weight at least
$n$. If $u_0=\wt L$, at most $u_0\lfloor n/q^b\rfloor$ positions
of $P_{n,c}L$ can meet $P_{n,c}H$: only kernel terms divisible by
$q^b$ can reach that part. Therefore
\[
 \wt(P_{n,c}D)\ge2n-2u_0\lfloor n/q^b\rfloor.
\]
Since $u_0\le u-1$ and $q^b\ge q$, this proves the bound. A strict
counterexample forces $q^b<2u_0\le2(u-1)$.
\end{proof}


\section{Residue bounds}
\label{sec:residue}

Group aggregate degrees modulo $p$. The governing estimate is simple:
if a correction has at least $N$ terms and the comparison polynomial
has $b$ terms in that class, the corrected class has at least $N-b$
terms and its weight changes by at least $N-2b$.
Lemma~\ref{lem:width-descent} supplies the lower bound $N$ for every
nonzero Frobenius correction. All weights below count monomials
in the individual fresh variables.

Let $K$ have characteristic $p>0$, let $0\ne C\in K[x_1,\ldots,x_a]$,
and put
\[
 c=\wt C,\qquad Z=z_1+\cdots+z_r,\qquad P=C+Z,\qquad r\ge1.
\]
For a base polynomial $B$, let $\sigma_p(B)$ be the largest number
of its support exponents congruent coordinatewise modulo $p$, with
$\sigma_p(0)=0$. Put $w_j=\binom{r+j-1}{j}$ for $0\le j<p$.
An expanded coefficient monomial multiplied by $Z^j$ costs $w_j$
fresh monomials. Replacing only this factor $Z^j$ by $(-C)^j$,
while holding all fresh $p$th powers fixed, contributes at most
$\sigma_p(C^j)$ positions to any fixed exponent class. These are
the cost and the possible cancellation that the reduction compares.

\paragraph{A residue-correction example.}
Over any field of characteristic two, take
\[
 C=1+x,\qquad Z=z_1+z_2,\qquad T=Z^2=z_1^2+z_2^2,
 \qquad P=C+Z.
\]
The multiple
\[
 G=P(P+T)=1+x^2+xT+ZT=A_0+ZA_1,
 \qquad A_0=1+x^2+xT,\quad A_1=T,
\]
has its aggregate degrees grouped modulo two. Replacing only the explicit
factor $Z$ by $-C=C$, while keeping the fresh square block $T$ fixed, gives
\[
 U=A_0+CA_1=1+x^2+T=P^2,\qquad V=G-U=PT.
\]
Since $ZT=z_1^3+z_1^2z_2+z_1z_2^2+z_2^3$, expansion in the individual
fresh variables gives $\wt G=\wt V=8$ and $\wt U=4$.
Aggregate degree modulo two is distinct from the full parity class
$(a,b_1,b_2)\bmod2$ of a monomial $x^az_1^{b_1}z_2^{b_2}$.
In the full zero class, $V$ contributes $T$, whereas $U$ contributes
$1+x^2+T$; their sum is $1+x^2$. Thus, with $b=2$, the corrected
class has exactly $\wt U-b=2$ terms, and its weight changes by
$\wt U-2b=0$. Extracting the square $U=P^2$ also returns $P$, whose
cofactor has base width zero, compared with base width one for the
cofactor $P+T$ of $G$. These identities illustrate the bookkeeping;
the minimality assumption below supplies the general lower bound $N$
for a nonzero correction class.

\begin{theorem}[Residue capacities]\label{thm:residue-descent}
Fix an integer $1\le N\le c+r$. If a nonzero $P$-multiple with fewer than $N$
terms exists, choose $G=PH$ as in Lemma~\ref{lem:width-descent} and write
\[
 G=\sum_{j=0}^{p-1}Z^jA_j(x,Z^p),\qquad m=\wt G<N.
\]
Define
\begin{align*}
 u_j&=\wt A_j,&
 R&=\sum_{j=1}^{p-1}w_ju_j,&
 M&=\sum_{j=1}^{p-1}\sigma_p(C^j)u_j,\\
 U&=\sum_{j=0}^{p-1}(-C)^jA_j,&
 V&=G-U.
\end{align*}
Then $m=u_0+R$, both $U$ and $V$ are $P$-multiples, and
$U=P^pE$ with $E\in K[x,Z^p]$. Every nonzero full exponent-residue
class of $U$ has at least $N$ terms.

If $h$ is the number of nonzero correction classes of $U$, then
\begin{equation}\label{eq:residue-capacity}
 m\ge\wt V+h(N-2M),\qquad h(N-M)\le m-R.
\end{equation}
More precisely, if $b_\alpha$ is the weight in class $\alpha$ of
the fresh-residue-zero part of $V$, then $b_\alpha\le M$ and
\[
 m\ge\wt V+\sum_{\alpha:\,U_\alpha\ne0}(N-2b_\alpha),
 \qquad
 \sum_{\alpha:\,U_\alpha\ne0}(N-b_\alpha)\le m-R.
\]
In particular, every class with $b_\alpha\le R$ has zero correction.

If $\kappa\le d(C)$ and $r\kappa\ge N$, the degree-zero aggregate
row of $G$ is nonzero, $u_0\ge\kappa$, and
\begin{equation}\label{eq:residue-budget}
 R\le N-1-\kappa.
\end{equation}
\end{theorem}

\begin{proof}
Every fresh exponent of $A_j$ is divisible by $p$. For $j<p$,
the monomials of $Z^j$ have exponents below $p$ and nonzero
multinomial coefficients. They occupy distinct fresh-residue
sectors, also distinct for different $j$. Hence $m=u_0+R$.

The polynomial
\[
 V=\sum_{j=1}^{p-1}\bigl(Z^j-(-C)^j\bigr)A_j
\]
is divisible by $P$. Also $U$ uses only fresh $p$th powers.
Evaluating $G$ at $Z=-C$, and writing $T=Z^p$, gives
$U(x,-C^p)=0$. Monic division in $K[x,T]$ yields
\[
 U=(T+C^p)E=P^pE,\qquad E\in K[x,Z^p].
\]
Every base exponent in $A_j$ occurs in $G$. Therefore, if $U\ne0$,
\[
 \operatorname{ord}_{x_i}U\ge\operatorname{ord}_{x_i}G,
 \qquad
 \deg_{x_i}U\le\deg_{x_i}G+(p-1)\deg_{x_i}C.
\]
Lemma~\ref{lem:width-descent} now gives at least $N$ terms in each
nonzero residue class $U_\alpha$.

The part of $V$ with zero fresh residues is
\[
 B=-\sum_{j=1}^{p-1}(-C)^jA_j.
\]
Each monomial of $A_j$ contributes at most $\sigma_p(C^j)$ terms
to a fixed full residue class, so $b_\alpha=\wt B_\alpha\le M$.
Adding a nonzero $U_\alpha$ leaves at least $N-b_\alpha$ terms and
changes weight by at least $N-2b_\alpha$. These classes are
disjoint and lie in $A_0$, proving the capacity inequalities.
If $b_\alpha\le R$, a nonzero correction would force
$u_0\ge N-b_\alpha\ge N-R>m-R=u_0$. Thus that class is protected.

Finally, write $G=\sum_{j\ge k}g_jZ^j$ with $g_k\ne0$.
Evaluation at $Z=-C$ and cancellation of $(-C)^k$ give
$g_k+\sum_{j>k}g_j(-C)^{j-k}=0$, so $C\mid g_k$.
If $k>0$, its row has at
least $W_r(k)\wt g_k\ge r\kappa\ge N$ terms. Hence $k=0$ and
$u_0\ge\wt g_0\ge\kappa$, proving \eqref{eq:residue-budget}.
\end{proof}

Three consequences will be used repeatedly. If $M\le R$, then
every correction class is protected and $U=0$: a nonzero correction
would leave at least $N-M$ terms in the zero fresh-residue part,
in addition to the $R$ terms already present elsewhere.
If $\wt V\ge N$
and $2M\le N$, equation~\eqref{eq:residue-capacity} excludes a
counterexample. If at most one correction class can occur and
$\wt V\ge\max\{N,2M\}$, it also excludes a counterexample.

\subsection{The characteristic-two form}

For $p=2$, write $D=A_1$. Then
\begin{equation}\label{eq:binary-residue-form}
 G=A_0+ZD,\qquad U=A_0+CD=P^2E,\qquad V=PD,
 \qquad m=\wt A_0+r\wt D.
\end{equation}
Here $A_0,D\in K[x,Z^2]$. The polynomial $D$ is nonzero: otherwise
$G=U$, and one of its nonzero correction classes would already
have at least $N$ terms. If $r\kappa\ge N$, then
\begin{equation}\label{eq:binary-u-budget}
 1\le u:=\wt D\le\left\lfloor\frac{N-1-\kappa}{r}\right\rfloor.
\end{equation}

The capacities can be sharpened using the parity distribution. Let
$s_\beta$ count the support of $C$ in base-parity class $\beta$,
and let $t_\gamma$ count the expanded terms of $D$ in base-parity
class $\gamma$. Then
\[
 b_\alpha\le\sum_\gamma t_\gamma s_{\alpha-\gamma}.
\]
If the positive $t_\gamma$ form the decreasing partition
$\lambda_1\ge\lambda_2\ge\cdots$ and the $s_\beta$ are sorted
decreasingly as $s_1\ge s_2\ge\cdots$, rearrangement gives the
uniform bound $b_\alpha\le\sum_i\lambda_i s_i$. Thus one may replace
$M$ by $\sum_i\lambda_i s_i$ in every capacity inequality.

If protection forces $U=0$, decompose
$D=\sum_\beta z^{2\beta}D_\beta(x)$. The products
$z^{2\beta}PD_\beta$ have disjoint supports because the fresh
exponents of $P$ are zero or one. Thus a short $PD$ would supply a
short base-only multiple.

\subsection{Finite-residue reduction}

\begin{corollary}[Uniform finite-residue formula]\label{cor:finite-residue}
If
\[
 \sigma_p(C^j)\le\binom{r+j-1}{j}\qquad(1\le j<p),
\]
then
\[
 d(P)=\min\left\{\sum_{j=0}^{p-1}w_j\wt B_j:
  B_j\in K[x],\ (B_0,\ldots,B_{p-1})\ne0,
  \sum_{j=0}^{p-1}(-C)^jB_j=0\right\}.
\]
In particular, a minimum multiple can be chosen with aggregate output
degree at most $p-1$ and cofactor degree at most $p-2$.
For $p=2$ and $r\ge\sigma_2(C)$ this becomes
\[
 d(C+Z)=\min_{D\ne0}\bigl(\wt(CD)+r\wt D\bigr).
\]
\end{corollary}
\begin{proof}
The stated inequalities give $M\le R$, so $U=0$ in
Theorem~\ref{thm:residue-descent}. Expand each $A_j$ in individual
fresh-$p$th-power monomials. For each such monomial, its coefficient
in $U=0$ says $\sum_j(-C)^jB_j=0$. Thus
$\sum_jZ^jB_j$ is a $P$-multiple of aggregate degree at most $p-1$.
These slices have disjoint expanded supports, so one is no heavier
than the original $G$. To obtain the minimum formula, if
$d(P)<c+r$ use the theorem with $N=d(P)+1$. If $d(P)=c+r$, the
polynomial $P$ already attains the minimum in the claimed form.
\end{proof}

\section{Intervals with \texorpdfstring{$n\ge67$}{n >= 67}}
\label{sec:large-weighted}

\begin{theorem}\label{thm:large-weighted}
Theorem~\ref{thm:field} holds for every $n\ge67$ and every field.
\end{theorem}

Throughout, remove the fresh primes in $(n/2,n]$ and rescale their
variables, writing
\[
 P_{n,c}=C+Z,\qquad r=\pi(n)-\pi(n/2),\qquad
 \wt C=n-r,
\]
and put $d_0=1+\pi(n/2)$ and $R=r(r+1)/2$. 

\subsection{Multipliers in the base variables}

For every nonzero $D$ in the base variables,
\begin{equation}\label{eq:large-base}
 \wt(CD)+r\wt D\ge n.
\end{equation}
Otherwise delete dispensable terms and normalize $D$; by the slice
argument of Section~\ref{sec:tools}, a deletion-minimal $D$ uses no
variable absent from $C$. It cannot be
monomial, which would give exactly $n$ terms. Its largest support
prime $q\le n/2$ satisfies $q<2(\wt D-1)$ by
Lemma~\ref{lem:directed-cut}, contradicting the localized estimate
(Theorem~\ref{arith:thm:localized-expansion}).

\subsection{A uniform affine reduction}

Lemma~\ref{arith:lem:no-core-fresh-budget} gives
$r\ge s:=\lfloor\sqrt n\rfloor\ge8$ for every $n\ge67$.
In characteristic $p>0$, support indices of $C^j$ are at most $n^j$.
A full exponent class modulo $p$ consists of integers $ab^p$ with
fixed $p$th-power-free part $a$, so
\[
 \sigma_p(C^j)\le\lfloor n^{j/p}\rfloor\qquad(1\le j<p).
\]
For $j=1$ this is at most $s\le r$. For $j=2<p$,
\[
 n^{2/p}\le n^{2/3}<(s+1)^{4/3}
 \le s(s+1)/2\le R,
\]
where the middle inequality follows from $s^3\ge8(s+1)$.
For $3\le j<p$,
\[
 \sigma_p(C^j)<n<(s+1)^2<\binom{s+2}{3}
 \le\binom{r+j-1}{j}.
\]
Corollary~\ref{cor:finite-residue} therefore gives a minimum multiple
of aggregate degree at most $p-1$. Since $P$ itself has $n$ terms,
a row of degree at least three
would already cost more than $n$ terms, so its degree is at most two.
In characteristic two it is at most one. In characteristic zero,
scalar compression and the same degree-three cost give degree at most
two directly. Monic division by $C+Z$ now shows that a minimum
multiple can be chosen with a base-only cofactor in characteristic two
and an affine cofactor $D+ZB$ otherwise. Characteristic two is settled
by \eqref{eq:large-base}.

In characteristic zero or odd characteristic, the exact expanded
weight of the affine product is
\begin{equation}\label{eq:large-affine}
 m=\wt(CD)+r\wt(D+CB)+R\wt B.
\end{equation}
If $B=0$, use \eqref{eq:large-base}. If $D=0$, the left side is
$r\wt(CB)+R\wt B$, which is at least
$\wt(CB)+r\wt B\ge n$. Otherwise $\wt(CD)\ge d_0$ by Lemma~\ref{lem:endpoints}.
For $n\ge229$, we have $d_0+R\ge n$. Indeed the thirteen blocks in
Appendix~\ref{sec:prime-blocks} prove this through $1799$.
For $1800\le n\le32768$, Lemma~\ref{arith:lem:no-core-fresh-budget}
gives $R>n^2/1800\ge n$; for $n\ge32768$ it gives $R>18n$.
Thus the affine weight is already at least $n$ throughout this range.

\subsection{Endpoint bounds for \texorpdfstring{$67\le n\le228$}{67 <= n <= 228}}

Put
\[
 \kappa=1+\pi(\sqrt n)+
       \sum_{\substack{\sqrt n<q\le n/2\\q\text{ prime}}}
          \left\lfloor\frac nq\right\rfloor,
 \qquad h=\pi(n/2)-\pi(\sqrt n).
\]
The core has the decomposition
\[
 C=A+\sum_{\substack{\sqrt n<q\le n/2\\q\text{ prime}}}x_qB_q,
\]
where $A$ is supported on the $\sqrt n$-smooth integers through $n$
and $B_q$ is a weighted interval of length $k_q=\lfloor n/q\rfloor\le15$.
No two of the displayed large prime variables occur together.
The initial-interval theorem gives $d(B_q)=k_q$, while the constant
and pure-prime vertices give $d(A)\ge1+\pi(\sqrt n)$.
Repeated endpoint slicing therefore gives
\[
 d(C)\ge\kappa.
\]
In characteristic different from two, the zero-large-variable row
and the pure $x_q^2$-rows of $C^2$ are $A^2$ and $B_q^2$.
They are disjoint. All these factors are nonmonomial, so the power-gap
lemma gives
\begin{equation}\label{eq:large-endpoint-square}
 \wt(C^2)\ge d(A)+1+\sum_q(d(B_q)+1)\ge\kappa+h+1.
\end{equation}

The endpoint table in Appendix~\ref{sec:prime-blocks} proves
\begin{equation}\label{eq:large-endpoint-budgets}
 \kappa+R+r\ge n,\qquad \kappa+R+h+1\ge n.
\end{equation}
Suppose the affine weight in \eqref{eq:large-affine} were less than
$n$, with $B,D\ne0$. Since $R\ge r$, the first inequality forces
$B$ to be monomial and then forces $D+CB=0$. The weight is consequently
$\wt(C^2)+R$, which is at least $n$ by
\eqref{eq:large-endpoint-square} and the second inequality.

This completes the proof of Theorem~\ref{thm:large-weighted}.

\begin{proof}[Proof of Theorem~\ref{thm:field}]
Theorem~\ref{thm:initial} proves the result for $1\le n\le24$.
Theorem~\ref{thm:finite} covers $25\le n\le66$. Theorem~\ref{thm:large-weighted} proves all
$n\ge67$. A nonzero scalar monomial multiplier has exactly $n$
output terms, proving sharpness. Laurentization preserves support
cardinality by Section~\ref{sec:tools}.
\end{proof}

\section{Coefficients and consequences}\label{sec:modules}

The field theorem says that erasing fewer than $n$ outputs preserves
injectivity. Over the ring of universal invertible weights, this gives
a left inverse. The same inverse proves the support bound for every
module, so erasure recovery and change of coefficients have a common
algebraic source.

\subsection{Erasure recovery}

Fix a finite nonempty input support $S\subset\N_{>0}$, let
$B=[n]S=\{ab:1\le a\le n,\ b\in S\}$, and erase
$E\subseteq B$ with $|E|<n$.
Over $A=\mathbb Z[t_1^{\pm1},\ldots,t_n^{\pm1}]$, form the
matrix $Q$ with rows $B\setminus E$, columns $S$, and entries
\[
 Q_{m,b}=\begin{cases}
 t_{m/b},&b\mid m\text{ and }1\le m/b\le n,\\
 0,&\text{otherwise}.
 \end{cases}
\]

\begin{proposition}[Universal erasure inverse]\label{prop:universal-inverse}
There is a matrix $L$ over $A$ with $LQ=I_{|S|}$. It remains
a left inverse when the $t_a$ are specialized to units in any
commutative ring.
\end{proposition}
\begin{proof}
Modulo every maximal ideal, $Q$ has full column rank by
Theorem~\ref{thm:field}: a nonzero kernel vector would have at
most $|E|<n$ outputs. In particular $Q$ has at least $|S|$ rows.
Thus the ideal of maximal minors lies
in no maximal ideal and equals $A$.
Choose $a_I\in A$ with $\sum_Ia_I\det Q_I=1$, where
$Q_I=E_IQ$ and $E_I$ selects $|S|$ rows. Then
\[
 L=\sum_Ia_I\operatorname{adj}(Q_I)E_I
 \quad\text{satisfies}\quad LQ=I_{|S|}.
\]
\end{proof}

\begin{remark}[Errors and erasures]\label{rem:errors-erasures}
The support bound also gives the standard unique-decoding criterion.
Suppose the output of weighted interval convolution is received with at
most $t$ corrupted coordinates outside a known set of $s$ erased
coordinates. The finitely supported input is uniquely determined whenever
$2t+s<n$. This holds over a field with nonzero weights and, by
Theorem~\ref{thm:module}, for module-valued inputs when every weight acts
injectively. Indeed, two compatible inputs $f,g$ have outputs differing
on at most $2t+s$ coordinates, whereas $f\ne g$ implies
$|\supp(T_c(f-g))|\ge n$. This is a uniqueness statement; it does not
assert an efficient decoding algorithm.
\end{remark}

\subsection{Modules and commuting actions}

\begin{theorem}[Module coefficients]\label{thm:module}
Let $M\ne0$ be a unital module over a commutative ring $R$,
and let $c_1,\ldots,c_n\in R$, where $n\ge1$. Every nonzero
finitely supported $f:\N_{>0}\to M$ satisfies
$|\supp(T_cf)|\ge n$ if and only if each $c_a$ acts injectively
on $M$. The bound is sharp.
\end{theorem}
\begin{proof}
Suppose first that the weights are units. If a nonzero input $f$
had fewer than $n$ outputs, put $S=\supp f$ and $E=\supp(T_cf)$.
In Proposition~\ref{prop:universal-inverse}, specialize $t_a=c_a$.
The surviving-output matrix $Q$ then acts on $M^S$ and annihilates
the vector of input values. Its left inverse gives $f=0$, a
contradiction.

For injective weight actions, localize at the multiplicative set
$U$ generated by the $c_a$. The map $M\to U^{-1}M$ is injective:
an element maps to zero only if some $u\in U$ kills it. Input and
output supports are unchanged, and all weights become units.

If some $c_a$ kills $m\ne0$, the singleton input $f(1)=m$ has
at most $n-1$ nonzero outputs. When all weight actions are injective,
every nonzero singleton input has exactly $n$ outputs.
\end{proof}

\begin{corollary}[Commuting group actions]\label{cor:commuting-actions}
Let $G\ne0$ be an abelian group and let
$\varphi_1,\ldots,\varphi_n$ be commuting endomorphisms. Then
\[
 (T_\varphi f)(m)=
 \sum_{\substack{1\le a\le n\\a\mid m}}\varphi_a(f(m/a))
\]
has at least $n$ nonzero values for every nonzero finitely supported
input if and only if every $\varphi_a$ is injective.
\end{corollary}
\begin{proof}
Apply Theorem~\ref{thm:module} to the commutative subring of
$\operatorname{End}_{\mathbb Z}(G)$ generated by the identity
and the $\varphi_a$.
\end{proof}

\subsection{Commutative monoids}

An endomorphism $\varphi$ of an additive monoid \emph{reflects zero}
if $\varphi(m)=0$ implies $m=0$. On a group this is injectivity;
on a monoid it is weaker, as $(x,y)\mapsto(x+y,0)$ on $\N^2$ shows.

\begin{theorem}[Monoid coefficients]\label{thm:monoid}
\label{thm:monoid-main}
Let $M\ne0$ be a commutative additive monoid and let
$\varphi_1,\ldots,\varphi_n$, $n\ge1$, be commuting
endomorphisms. Every nonzero finitely supported input satisfies
$|\supp(T_\varphi f)|\ge n$ if and only if each $\varphi_a$
reflects zero. The bound is sharp.
\end{theorem}
\begin{proof}
Let $M^\times$ be the group of additively invertible elements of $M$,
and suppose each weight reflects zero. Every endomorphism preserves
$M^\times$, and its restriction to this group is injective. Put
$\psi=\varphi_1\cdots\varphi_n$; this map also reflects zero.

If $\psi(f(b))\notin M^\times$ for some $b$, then
$\varphi_a(f(b))\notin M^\times$ for every $a$: applying the remaining
endomorphisms to a unit would give a unit. Each output at $ab$,
$1\le a\le n$, therefore contains a nonunit summand and is nonzero,
since a summand of a zero sum has an additive inverse.

Otherwise $h=\psi\circ f$ is a nonzero input with values in
$M^\times$. Apply Corollary~\ref{cor:commuting-actions} to the
restrictions of the weights to this group. Commutativity gives
$T_\varphi h=\psi\circ T_\varphi f$, and the zero-reflecting property
of $\psi$ makes these two outputs have identical supports.
Thus the required bound follows in this case as well.

A weight that kills a nonzero element gives a singleton input with
fewer than $n$ outputs. Zero-reflecting weights give exactly $n$
outputs on every nonzero singleton input.
\end{proof}

The theorem also applies to semimodules over commutative semirings.
Commutativity of the weight actions is essential. On $\mathbb Q^2$, put
\[
 B=\begin{pmatrix}1&0\\0&-1\end{pmatrix},\qquad
 C=\begin{pmatrix}0&1\\1&0\end{pmatrix}.
\]
Since $B^2=C^2=I$ and $BC+CB=0$,
\[
 (I+Bx+Cy+\tfrac12Ix^2)(I-Bx-Cy+\tfrac12Ix^2)
 =I+\tfrac14Ix^4-Iy^2.
\]
With $x=X(2)$ and $y=X(3)$, applying this identity to a nonzero
vector gives four invertible interval weights and only three outputs.

\subsection{Rank and equality}

\begin{theorem}[Coefficient rank]\label{thm:coefficient-rank}
Let $n\ge1$, let $V$ be a vector space over $K$, let $c_1,\ldots,c_n\in K^\times$,
and let $f:\N_{>0}\to V$ be nonzero and finitely supported. If
$\rho=\dim_K\operatorname{span}\{f(b):b\ge1\}$, then
\[
 |\supp(T_cf)|\ge n+\rho-1.
\]
\end{theorem}
\begin{proof}
The input values and the output values $g=T_cf$ span the same space.
One inclusion is immediate; the other follows by induction on $b$ from
\[
 c_1f(b)=g(b)-
 \sum_{\substack{2\le a\le n\\a\mid b}}c_af(b/a).
\]
Choose $\rho-1$ independent output values and a nonzero linear
functional on $\Span\{f(b):b\ge1\}$ vanishing on them. The projected
input is nonzero, so Theorem~\ref{thm:field} leaves at least $n$
output positions. At least $\rho-1$ original positions were removed.
\end{proof}

\begin{corollary}[Equality at a prime index]\label{cor:prime-rigidity}
\label{cor:vector-prime-rigidity}
Let $V$ be a vector space over $K$, let $\ell\ge3$ be prime,
let $c_1,\ldots,c_\ell\in K^\times$, and put $P=P_{\ell,c}$.
Every vector-valued Laurent multiple
$G=PH$ with exactly $\ell$ nonzero coefficients has the form
\[
 G=vMP^q,\qquad H=vMP^{q-1},
\]
where $v\in V\setminus\{0\}$, $M$ is a Laurent monomial, and
$q=1$ in characteristic zero or a power of the characteristic
otherwise. If $G$ is a polynomial, $M$ has nonnegative exponents.
\end{corollary}
\begin{proof}
A monomial translation makes $H$ polynomial. The rank bound forces
its coefficients to span a line, so $H=vh$ for a scalar polynomial $h$.
By Theorem~\ref{thm:field},
$d(P_{\ell-1,c})=\wt P_{\ell-1,c}=\ell-1$.
Set $z=c_\ell x_\ell$, a fresh variable. The strict case of
Theorem~\ref{thm:completion} gives $Ph=aMP^q$; absorb $a$ into $v$
and undo the initial monomial translation.
Finally, $P^q$ has nonzero constant term, so $MP^q$ is polynomial
only when $M$ has no negative exponent.
\end{proof}

\subsection{Consequences for the original Pilz conjecture}

For a finite nonempty set $A\subset\N_{>0}$, write
\[
 S_n(A)=\mathop{\triangle}_{a\in A}\{a,2a,\ldots,na\},
 \qquad F_A=\sum_{a\in A}X(a)\in\mathbb F_2[\mathbf x],
 \qquad P_n=\sum_{a=1}^n X(a).
\]

\begin{corollary}[All equality cases at prime indices]
\label{cor:binary-prime-equality}
Let $\ell\ge3$ be prime. Then $|S_\ell(A)|=\ell$ if and only if
there exist $b\ge1$ and $e\ge0$ such that
\[
 F_A=X(b)P_\ell^{\,2^e-1}.
\]
For these same $b,e$, the output is exactly
\[
 S_\ell(A)=\{ba^{2^e}:1\le a\le\ell\}.
\]
\end{corollary}
\begin{proof}
Apply Corollary~\ref{cor:prime-rigidity} over $\mathbb F_2$ to
$P_\ell F_A$. Its nonzero scalar is $1$, so
$F_A=MP_\ell^{2^e-1}$ for a monomial $M$ with nonnegative exponents.
The constant term of $P_\ell^{2^e-1}$ is $1$, including when $e=0$.
Thus $M$ occurs in $F_A$, and $M=X(b)$ for some $b\in A$.
Conversely, the displayed form gives
\[
 P_\ell F_A=X(b)P_\ell^{2^e}
 =\sum_{a=1}^{\ell}X(ba^{2^e})
\]
by Frobenius. These $\ell$ monomials are distinct, proving equality
and the claimed output set.
\end{proof}

The case $e=0$ gives the singleton inputs $A=\{b\}$. The case
$e=1$ gives $A=b[\ell]$ and output $\{ba^2:1\le a\le\ell\}$;
higher $e$ give the remaining equality families through the same
polynomial formula.

\begin{corollary}[Sharp parity refinement]\label{cor:binary-parity}
For every $n\ge1$,
\[
 |S_n(A)|\equiv n|A|\pmod2.
\]
Consequently, if $n$ is odd and $|A|$ is even, then
$|S_n(A)|\ge n+1$. This bound is sharp for every odd $n$.
\end{corollary}
\begin{proof}
Counting the $n|A|$ multiplicative representations modulo two gives
the parity identity, also recorded in Holdridge--Pach~\cite{HP}.
Theorem~\ref{thm:field} gives $|S_n(A)|\ge n$; an even integer at
least the odd integer $n$ is at least $n+1$. For $A=\{1,2\}$,
the intersection $[n]\cap2[n]$ has $\lfloor n/2\rfloor$ elements, so
\[
 |S_n(\{1,2\})|=2n-2\lfloor n/2\rfloor=n+1
 \qquad(n\text{ odd}).\qedhere
\]
\end{proof}

\appendix

\section{The initial intervals through twenty-four}
\label{sec:initial}

\begin{theorem}\label{thm:initial}
For every field $K$, every $1\le n\le24$, and every choice
$c_a\in K^\times$, the polynomial
\[
 P=\sum_{a=1}^n c_aX(a)
\]
satisfies $d(P)=n$.
\end{theorem}

Throughout this section, all coefficients in a displayed support are
nonzero. Different occurrences of that support may have different
coefficients. A dense polynomial of degree $e$ has all $e+1$
coefficients nonzero. We use $x=x_2$, $y=x_3$, $z=x_5$, $w=x_7$,
and $v=x_{11}$. Coefficients of fresh variables are absorbed by rescaling.
An exponent is \emph{distinguished} if it is $1$ in characteristic zero
or a power of the characteristic otherwise. Distinguished powers preserve
weight.

\subsection{Three elementary estimates}

The endpoint bound says that if the lowest and highest coefficients in
one variable are $A$ and $B$, every nonzero multiple has at least
$d(A)+d(B)$ terms. Each additional occupied row contributes at least one
more term. We also use that a nonmonomial polynomial has no monomial
multiple.

\begin{lemma}[A quadratic obstruction]\label{lem:initial-quadratic}
Let $A,B,D\in K[x]$ be nonzero.  If
\[
 \frac{B^2}{AD}
\]
is nonconstant, every nonzero multiple of $A+By+Dy^2$ has at least three
occupied $y$-rows.
\end{lemma}
\begin{proof}
Suppose there were at most two occupied rows. After removing the
lowest power of $y$, the quadratic would divide $U+Vy^k$ over $K(x)$,
where $UV\ne0$ and $k\ge1$. Put
\[
 s=\frac{AD}{B^2},\qquad y=-\frac BD t.
\]
The quadratic becomes a nonzero scalar multiple of $t^2-t+s$.
Modulo this monic quadratic, write
\[
 t^k=a_k(s)t+b_k(s).
\]
The coefficients are polynomials in $s$ over the prime field, with
$a_1=a_2=1$ and $a_{k+1}=a_k-sa_{k-1}$. In particular every $a_k$
has constant coefficient $1$, so none is the zero polynomial.
The binomial divisibility would force $a_k(s)=0$. A nonconstant
rational function in $x$ is transcendental over $K$, so this is
impossible. Indeed, after writing such a function as a reduced
fraction, clearing denominators in a polynomial relation first
forces its denominator to be constant, and then comparison of
degrees forces its numerator to be constant as well.
\end{proof}

\begin{lemma}[Two rows and a binomial coefficient]
\label{lem:initial-extraction}
Let $R\in K[x,\mathbf y]$ and let $D=d_0+d_1x$, with $d_0d_1\ne0$.
Suppose $\gcd(R,D)=1$.
If $R+Dt$ divides $U+Vt^k$, where $U,V\in K[x,\mathbf y]$,
$UV\ne0$, and $k\geq1$, then, up to signs,
\[
 U=R^kM,\qquad V=D^kM
\]
for the same polynomial $M\in K[x,\mathbf y]$.
If $\wt(V)\leq3$, then $M$ is a monomial in $\mathbf y$ times a
univariate polynomial in $x$.

In particular, if $d(R)\geq a$ and
\[
 \wt\bigl(R^kM_0(x)\bigr)>a
 \quad\text{for all }k\geq1\text{ and }0\ne M_0\in K[x],
\]
then $d(R+Dt)\geq a+3$.
\end{lemma}
\begin{proof}
Substitution $t=-R/D$ and unique factorization give the first assertion.
Each occupied $\mathbf y$-row of $D^kM$ is divisible by the nonmonomial
$D^k$, so it has at least two terms.  A polynomial of weight at most three
can therefore have only one such row.

For the last assertion, endpoints give $a+2$.  A multiple with at most
$a+2$ terms must have exactly two $t$-rows, with a binomial upper row and
a lower row of weight at most $a$.  Extraction makes the lower row a
monomial times $R^kM_0(x)$, contrary to the assumed bound.
\end{proof}

\begin{lemma}[Power estimates]\label{lem:initial-powers}
Every positive power of a dense univariate polynomial of degree at least
two has at least three terms.  Consider also
\[
 A=Q(x)+B(x)y+D(x)y^2,
\]
where $Q$ is dense of degree at least three and $B$ is dense of degree at
least one.  In either of the following situations,
\[
 \wt(A^kM(x))\geq6\qquad(k\geq1,\ M\ne0):
\]
\begin{enumerate}
 \item $D$ is a nonzero constant;
 \item $D$ is a binomial with nonzero constant and linear coefficients.
\end{enumerate}
In the first situation $\wt(A^k)\geq\deg(B)+5$; in the second,
$\wt(A^k)\geq7$.

Let $L=l_0+l_1x+l_2y$ and
$T=t_0+t_1x+t_2x^2+t_3y$, with all coefficients nonzero.  Then
\[
 d(L)=3,\qquad d(T)=4,\qquad
 \wt(L^k)\geq3,\qquad\wt(T^k)\geq4.
\]
If $M(x)$ is nonmonomial, both $L^kM$ and $T^kM$ have at least four
terms.  A nondistinguished power of $L$ has at least six terms, and a
nondistinguished power of a binomial has at least three.
\end{lemma}
\begin{proof}
First let the exponent $h$ be nonzero in the field.  For a dense
univariate polynomial of degree $e\geq2$, the constant, linear, and
leading coefficients of its $h$th power are nonzero, at the distinct
degrees $0,1,eh$.

The $y$-rows of $A^hM$ at degrees $2h,2h-1,0$ are
\[
 D^hM,\qquad hD^{h-1}BM,\qquad Q^hM.
\]
If $D$ is constant and $M$ is nonmonomial, each row has at least two
terms.  If $M$ is a monomial, their weights are at least
$1,\deg(B)+1,3$.  If $D$ is a binomial, all three rows have at least two
terms; when $M$ is a monomial, the bottom one has at least three.

In characteristic $p$, write $k=qh$ with $q$ a power of $p$
and $p\nmid h$. Split $A^{qh}M$ into the residue classes of the
$x$-exponent modulo $q$. Applying Lemma~\ref{lem:frobenius-extraction}
to any nonzero class gives a support-contained polynomial
$x^i(A^hM_i(x))^q$ with $M_i\ne0$.
The bounds just proved for $h$ therefore imply the bound six for
arbitrary $k$ and $M$. For monomial $M$, the sharper power estimates
follow directly from Frobenius. The univariate assertion follows
the same way, or simply by applying Frobenius to the reduced power.

The distance of $L$ follows from endpoints.  The distance of $T$ follows
from sharp completion, Theorem~\ref{thm:completion}, applied to a dense
quadratic in $x$: $\min(3+1,2+2)=4$.  The asserted bounds with nonmonomial
$M$ follow from the two $y$-endpoints, each of which has at least two
terms.  The bound for $T^k$ follows from its monomial top row and the
three-term bound for the bottom dense-quadratic power.  Finally, the
ordinary multinomial formula and its base-$p$ digit version give at least
six terms for a nondistinguished power of $L$, and at least three for a
nondistinguished binomial power.
\end{proof}

\subsection{A common test for adjoining fresh variables}

Two support observations will supply the geometric part of the test.
First, if a support meets at least $b$ lines in every nonzero direction,
a binomial multiple has at least $2b$ terms: the two endpoints on every
occupied line survive. This remains true after translation or Frobenius
dilation.

Second, consider
\begin{equation}\label{eq:initial-sliced-support}
 C=A(x,y)+\sum_{i=1}^s z_iB_i(x,y).
\end{equation}
Its projection to the $z_i$-coordinates is the simplex
$\{0,e_1,\ldots,e_s\}$. A translate of this simplex and a translate
of its $q$-dilate, $q\ge2$, meet in at most one point: two intersections
would equate a nonzero vector with coordinates in $\{-1,0,1\}$ to $q$
times another such vector. Thus two translated unequal distinguished
powers of $C$ overlap in at most one projected row. An $A$-to-$A$
intersection has at most as many points as the largest exponent class
of $A$ modulo $q$; any other intersection has at most the smaller
coefficient support's size.

Write $W_r(j)=\wt((u_1+\cdots+u_r)^j)$. The digit formula gives
$W_r(j)\ge r$ for $j>0$. For $r\ge2$, equality occurs exactly at
distinguished degrees. For $r=2,3,4$, every nondistinguished degree
costs at least $3,6,10$, respectively. For $r=2,3$, equality in
this lower bound occurs only in characteristic zero or odd
characteristic, at twice a distinguished degree.

\begin{lemma}[A short fresh-row test]\label{lem:short-fresh-test}
Let $P=C+Z$, where $Z$ is a sum of $r\ge2$ fresh variables.
Put $c=\wt C$, $N=c+r$, and suppose $d(C)\ge\kappa$,
$N\le r(\kappa+1)$ and $N-1-\kappa\le2r$.
To prove $d(P)=N$, it suffices to check:
\begin{enumerate}
 \item $\wt(C^j)+W_r(j)\ge N$ for every nondistinguished degree
 $j$ with $W_r(j)\le N-1-\kappa$;
 \item every binomial multiple of a distinguished power of $C$
 has at least $N-2r$ terms;
 \item every nonzero sum of monomial translates of two unequal
 distinguished powers of $C$ has at least $N-2r$ terms.
\end{enumerate}
For $r=2$ the same conclusion holds when $N-1-\kappa=5$,
provided one also checks
\[
 \wt(MC^q+LC^{2q'})\ge N-5
\]
for nonzero scalar monomials $M,L$ and distinguished $q,q'$
in characteristic zero or odd characteristic.
\end{lemma}
\begin{proof}
Compress a hypothetical multiple of weight at most $N-1$ and
write it as $\sum_j A_jZ^j$. The coefficient of its lowest occupied
$Z$-degree is divisible by $C$. If that degree were positive, the row would
cost at least $r\kappa$. There must be another occupied row,
since a single nonzero row cannot vanish at $Z=-C$; it costs
at least $r$. This contradicts $N\le r(\kappa+1)$.

Thus $A_0\ne0$, $\wt A_0\ge\kappa$, and the positive rows
cost at most $2r$. Every such row costs at least $r$, and a
nondistinguished degree costs strictly more than $r$.
There are only three possibilities: a single monomial
coefficient, a single binomial coefficient at a distinguished
degree, or two monomial coefficients at distinct distinguished
degrees. Evaluation at $Z=-C$ gives exactly the three
conditions in the statement. A single distinguished monomial
row already costs $c+r=N$.

If $r=2$ and the budget is five, the only extra possibility is
two monomial rows of costs two and three. A cost of three is
possible only in characteristic zero or odd characteristic,
at a degree $2q'$ with $q'$ distinguished. The additional
condition excludes it.
\end{proof}

\subsection{The intervals through eight}

The cases $n=1,2,3$ are, respectively, a monomial, a binomial, and a
three-term affine simplex; their distances are $1,2,3$.  At $n=4$, the
polynomial is a dense quadratic in $x$ plus a fresh variable.  Sharp
completion gives $\min(3+1,2+2)=4$.  Adding the prime-$5$ variable gives
$n=5$.

For $n=6$ and $n=7$, remove the one or two fresh prime terms.  The
remaining five-term core is
\[
 Q_2(x)+D(x)y,
\]
where $Q_2$ is dense quadratic and $D$ is a nondegenerate binomial.
Its two endpoint rows each have distance at least two, so its distance
is at least four.  Completion gives
$\min(5+r,4+2r)=5+r$ for $r=1,2$.
At $n=8$, replace $Q_2$ by a dense cubic.  The core now has six terms
and still has distance at least four; the two fresh variables give
$\min(6+2,4+4)=8$.

\subsection{Five quadratic bases}

The core estimates all start from the following polynomials:
\[
 A_\nu=Q_a(x)+B_b(x)y+D(x)y^2.
\]
Here $Q_a,B_b$ are dense of the indicated degrees, and $D$ is either a
nonzero constant or a nondegenerate binomial.
\begin{center}
\begin{tabular}{c|rrl|rrr}
\toprule
$\nu$&$a$&$b$&$D$&$\wt(A_\nu)$&$d(A_\nu)\ge$&$\wt(A_\nu^k)\ge$\\\midrule
9&3&1&constant&7&4&6\\
12&3&2&constant&8&4&7\\
16&4&2&constant&9&4&7\\
18&4&2&binomial&10&5&7\\
24&4&3&binomial&11&5&7\\
\bottomrule
\end{tabular}
\end{center}
For all five bases, Lemma~\ref{lem:initial-powers} gives
$\wt(A_\nu^kM(x))\ge6$ whenever $k\ge1$ and $M\ne0$.
Except at $\nu=16$, the numerator and denominator of $B_b^2/(Q_aD)$
have different degrees. Lemma~\ref{lem:initial-quadratic} therefore
forces three $y$-rows. Their endpoint costs give the distance entries.
In particular,
\begin{align}
 d(A_9)&\ge4,&\wt(A_9^k)&\ge6,\label{eq:initial-nine-core}\\
 d(A_{12})&\ge4,&\wt(A_{12}^kM)&\ge6,&
 \wt(A_{12}^k)&\ge7,\label{eq:initial-twelve-A}\\
 d(A_\nu)&\ge5,&\wt(A_\nu^kM)&\ge6,&
 \wt(A_\nu^k)&\ge7\quad(\nu=18,24).
 \label{eq:initial-late-A}
\end{align}

At $\nu=16$, only the case in which $B_2^2/(Q_4D)$ is constant
needs attention.
Write $D=g$ and $Q_4=\lambda B_2^2$, with $g\lambda\ne0$.
A three-term multiple would have two $y$-rows: a binomial lower row
$U(x)$ and a monomial upper row $V(x)y^k$. Substitute $y=B_2t$ and
reduce modulo $\lambda+t+gt^2$. In this two-dimensional quotient
$t$ is a unit, so the remainder $a+bt$ of $t^k$ is nonzero.
Divisibility forces $b=0$, hence $a\ne0$ and $U=-aVB_2^k$.
Lemma~\ref{lem:initial-powers} gives at least three terms in $B_2^k$,
a contradiction. Thus $d(A_{16})\ge4$ as well.

\subsection{The core extensions}

Let $L=l_0+l_1x+l_2y$, let
$T=t_0+t_1x+t_2x^2+t_3y$, and let $D_i$ denote nondegenerate
binomials in $x$. The following table describes the cores obtained by
removing the primes in $(n/2,n]$ from the indicated intervals.
The last column is a lower bound for their distance.
\begin{center}
\begin{tabular}{r|l|rr}
\toprule
$n$&$C_n$&$\wt(C_n)$&$\kappa$\\\midrule
9&$A_9$&7&4\\
12&$A_{12}+D_1z$&10&7\\
14&$A_{12}+D_1z+D_2w$&12&10\\
15&$A_{12}+Lz+D_2w$&13&10\\
16&$A_{16}+Lz+D_2w$&14&10\\
18&$A_{18}+Lz+D_2w$&15&11\\
20&$A_{18}+Tz+D_2w$&16&12\\
21&$A_{18}+Tz+Lw$&17&12\\
22&$A_{18}+Tz+Lw+D_3v$&19&15\\
24&$A_{24}+Tz+Lw+D_3v$&20&15\\
\bottomrule
\end{tabular}
\end{center}
We prove the entries by three applications of the same endpoint argument.

First, put $E=A_{12}+D_1z$. The constant $y^2$ coefficient gives
$\gcd(A_{12},D_1)=1$. Lemma~\ref{lem:initial-extraction}, followed by
the power endpoint bounds, gives
\begin{equation}\label{eq:initial-twelve-E}
 d(E)\ge7,\qquad \wt(E^kM(x))\ge8,\qquad \wt(E^k)\ge9.
\end{equation}
A second binomial extension therefore has distance at least $10$,
proving the entries for $C_{12}$ and $C_{14}$.

Next put $F_\nu=A_\nu+Lz$ for $\nu=12,16,18,24$.
Endpoints give $d(F_\nu)\ge d(A_\nu)+3$, and
\begin{equation}\label{eq:initial-fifteen-F}
 \wt(F_\nu^kM(x))\ge10\qquad(k\ge1, M\ne0).
\end{equation}
Indeed, for nonmonomial $M$ the two endpoint costs are at least $6+4$;
for monomial $M$ they are at least $7+3$. In particular,
\begin{equation}\label{eq:initial-late-F}
 d(F_\nu)\ge8,\qquad \wt(F_\nu^kM(x))\ge10
 \quad(\nu=18,24).
\end{equation}
The coefficient of $yz$ in $F_\nu$ is a nonzero constant, so
$\gcd(F_\nu,D_2)=1$. A binomial extension raises the stated distance
bound by three, giving the entries for $C_{15},C_{16},C_{18}$.
For $C_{15}$ and $C_{16}$ we also need a bare-power bound:
\[
 \wt(C_n^k)\ge16\qquad(n=15,16; k\text{ nondistinguished}).
\]
Remove the characteristic-power factor of $k$, leaving $h\ge2$
nonzero in the field. The $z$-endpoints of $F_\nu^h$ cost at least
$7+6=13$, and the binomial $w$-endpoint of $C_n^h$ costs at least
three. Frobenius preserves this bound.

Finally, for $\nu=18,24$ put $E_\nu=A_\nu+Tz$ and
$H_\nu=E_\nu+Lw$. Endpoints give
\begin{equation}\label{eq:initial-late-E}
 d(E_\nu)\ge9,\qquad
 \wt(E_\nu^kM(x))\ge10,\qquad \wt(E_\nu^k)\ge11.
\end{equation}
Here the power costs are $6+4$ for nonmonomial $M$ and $7+4$ for
monomial $M$. Consequently $d(H_\nu)\ge12$, and
\begin{equation}\label{eq:initial-twentyone-power}
 \wt(H_\nu^kM(x))\ge14.
\end{equation}
The latter costs are $10+4$ or $11+3$, according as $M$ is
nonmonomial or monomial. The coefficients of $yz$ in $E_\nu$ and
of $yw$ in $H_\nu$ ensure coprimality with every $D_i$.
Lemma~\ref{lem:initial-extraction} now gives distance at least $12$
for $E_{18}+D_2w$, and at least $15$ for $H_\nu+D_3v$.
This proves the remaining entries.

\subsection{The mixed power at sixteen}

Only one additional mixed-power estimate is needed. Write
$C_{16}=A_{16}+Lz+Dw$, as in the core table. In characteristic
different from two, its square has six distinct $(z,w)$-rows,
with coefficients
\[
 A_{16}^2,\quad 2A_{16}L,\quad 2A_{16}D,\quad L^2,\quad 2LD,\quad D^2.
\]
Their weights are at least $4,4,4,4,4,3$, respectively. The first
three bounds follow from $d(A_{16})\ge4$. The power gap gives
$\wt(L^2)\ge4$; the two $y$-endpoints of $LD$ each cost two;
and a nondegenerate binomial square has three terms.

For distinguished $q,q'$, these six rows remain distinct in
$C_{16}^{2q'}$, while a monomial translate of $C_{16}^q$ occupies
only three $(z,w)$-rows. At least three square rows are untouched.
Thus, for nonzero scalar monomials $M,N$,
\begin{equation}\label{eq:initial-sixteen-mixed}
 \wt(MC_{16}^q+NC_{16}^{2q'})\ge4+4+3=11.
\end{equation}

\subsection{Applying the fresh-row test}

We now collect the estimates needed in Lemma~\ref{lem:short-fresh-test}.
In the table, $b$ is a lower bound for the number of parallel lines in
any direction, and $o$ bounds the overlap of two translated unequal
distinguished dilates of the core support.
\begin{center}\small
\begin{tabular}{r|rrrrr|l}
\toprule
$n$&$c$&$r$&$\kappa$&$b$&$o$&
relevant nondistinguished powers\\\midrule
9&7&2&4&3&3&weight at least $6$\\
12&10&2&7&4&3&weight at least $9$\\
15&13&2&10&6&3&weight at least $16$\\
16&14&2&10&6&4&weight at least $16$; mixed bound $11$\\
18&15&3&11&6&4&weight at least $12$\\
21&17&4&12&7&4&none fit the budget\\
22&19&3&15&8&4&weight at least $16$\\
24&20&4&15&8&4&none fit the budget\\
\bottomrule
\end{tabular}
\end{center}
Here is a uniform verification of the geometric columns.
Each base $A_\nu$ requires at least three lines: a horizontal direction
sees three $y$-levels, and any other direction separates its four or five
$x$-axis points. Each $L$- or $T$-slice needs at least two lines, and
each $D_i$-slice at least one. Adding these counts gives $b$ for
directions in the $xy$-plane. A direction outside that plane separates
all the points of $A_\nu$, giving at least the same bound.

For overlaps, use~\eqref{eq:initial-sliced-support}. Every exponent
class of a cubic base modulo $q\ge2$ has at most three points, and
every class of a quartic base has at most four. For $q=2$, the largest
possible quartic class is $\{1,x^2,x^4,y^2\}$; for $q\ge3$,
there are at most two points in a class. The other coefficient slices
have at most three terms in the cubic cases and four in the quartic
cases. This proves the column $o$, including the unsliced base $C_9$.

In every row, $n\le r(\kappa+1)$ and the positive-row budget
$n-1-\kappa$ is at most $2r$, except at $n=16$, where it is five.
The binomial and unequal-dilate conditions follow from
\[
 2b\ge n-2r,\qquad 2c-2o\ge n-2r.
\]
The power estimates prove the remaining condition. At $n=18,22$,
the power gap gives weights at least $\kappa+1=12,16$, respectively;
adding the fresh-row cost of at least six gives $n$. At $n=21,24$, every nondistinguished degree
costs at least ten, exceeding the budget eight.
At $n=16$, use~\eqref{eq:initial-sixteen-mixed} as well.
Thus all eight rows follow from Lemma~\ref{lem:short-fresh-test}.

The two cores omitted from this last table follow directly from sharp
completion:
\[
 d(P_{14})\ge\min(12+2,10+4)=14,\qquad
 d(P_{20})\ge\min(16+4,12+8)=20.
\]
Finally, adjoining the fresh prime variable gives the cases
$13,17,19,23$ from $12,16,18,22$, respectively. Together with the
cases through eight, this leaves only $10,11$.

\subsection{The tenth and eleventh intervals}

Let $C=A_9$, let $D=d_0+d_1x$, and set
\[
 E=C+Dz,\qquad P=E+w.
\]
The support of $P$ is the tenth interval.  We first establish the
following statements about its nine-term core $E$:
\begin{align}
 &d(E)\geq7,\label{eq:initial-ten-distance}\\
 &\text{every $E$-multiple of weight at most seven has three $z$-rows
 of weights $(4,1,2)$;}\label{eq:initial-ten-shape}\\
 &\wt(E^j)\geq9\quad(j\geq1),\label{eq:initial-ten-power}\\
 &\wt(HE^j)\geq8\quad\text{if $H$ is a binomial.}
 \label{eq:initial-ten-binomial}
\end{align}

The strengthened core estimate
\[
 \wt(C^jM(x))\geq6\qquad(M\ne0)
\]
is Lemma~\ref{lem:initial-powers}.  Endpoints first give $d(E)\geq6$.
A multiple of weight at most seven with exactly two $z$-rows has lower
row of weight at least four and upper row of weight at most three.
Since the coefficient of $y^2$ in $C$ is a nonzero constant,
$\gcd(C,D)=1$.  Lemma~\ref{lem:initial-extraction} makes the lower row
a monomial times $C^jM(x)$, costing at least six, and the upper row
costs at least two.  This is impossible.  There are therefore at least
three rows.  Their endpoint costs are four and two, which proves
\eqref{eq:initial-ten-distance} and the forced shape
\eqref{eq:initial-ten-shape}.

Distinguished powers of $E$ have nine terms.  For any other exponent,
remove its largest characteristic-power factor and use a reduced
exponent $h\geq2$ that is nonzero in the field.  The $z$-rows of $E^h$
at degrees $h,h-1,0$ are
\[
 D^h,\qquad hD^{h-1}C,\qquad C^h.
\]
Their weights are at least $2,4,6$, giving at least twelve terms and
proving~\eqref{eq:initial-ten-power}.

For~\eqref{eq:initial-ten-binomial}, write $j=qh$, with $q$
distinguished and $h$ nonzero in the field; in characteristic zero take
$q=1$.  The support of $E^j$ contains the three $x$-axis points at
exponents $0,q,3j$.  Its support also contains the four affinely independent endpoint
points $(0,0,0),(3j,0,0),(0,2j,0),(0,0,j)$, so it has affine
dimension three.  In a direction
not parallel to the $x$-axis these three points lie on different
parallel lines.  If those were the only occupied lines, the support
would lie in the plane generated by that direction and the $x$-axis,
a contradiction.  In the $x$-direction, four nonzero rows have
$(y,z)$-coordinates
\[
 (0,0),\quad(2j,0),\quad(2j-q,0),\quad(0,j).
\]
They come from the bottom, top, next-to-top, and final endpoint rows
described above.  Thus every direction meets at least four lines,
and a binomial product has at least eight terms.

We need one last consequence of the forced shape:
\begin{equation}\label{eq:initial-ten-two-powers}
 \wt(ME^i+NE^j)>7
 \quad(1\leq i<j)
\end{equation}
for any nonzero scalar monomials $M,N$, provided the sum is nonzero.
Suppose otherwise.  Write $M=z^aM_0(x,y)$ and $N=z^bN_0(x,y)$.
If $a\ne b$, the lowest $z$-row is a monomial times a power of $C$,
which costs at least six, contradicting its forced weight four.
Thus $a=b$, and we remove this common $z$-power.  The resulting
constant row is nonzero: its vanishing would make the positive power
$C^{j-i}$ a scalar Laurent monomial, impossible for a nonmonomial
$C$.  The top row is $N_0D^j$ and has exactly two terms.  Hence $j$ is
distinguished.  Characteristic zero is impossible because $j>i\geq1$.
In characteristic $p$, $E^j$ has just the two $z$-rows of degrees zero
and $j$.  Every positive $z$-row of $M_0E^i$ below degree $j$ survives
unchanged.  Exactly three rows in the sum force $E^i$ to have only two
rows, so the binomial digit formula forces $i$ to be distinguished as
well.  The middle row is then $M_0D^i$, of weight two, contradicting
its forced weight one.  This proves~\eqref{eq:initial-ten-two-powers}.

Now suppose a $P$-multiple has at most nine terms.  Remove its lowest
$w$-power, an operation that preserves both weight and divisibility,
and write it as
\[
 U+\sum_{j>0}A_jw^j,\qquad U\ne0.
\]
Substitution $w=-E$ shows $E\mid U$, so $\wt(U)\geq7$.  At most two
terms remain at positive $w$-degrees.  If these form one monomial row,
\eqref{eq:initial-ten-power} gives total weight at least $9+1=10$.
If they form one binomial row, \eqref{eq:initial-ten-binomial} gives
at least $8+2=10$.  If they form two monomial rows,
\eqref{eq:initial-ten-two-powers} excludes the required constant row
of weight at most seven.  This proves $n=10$.  Adding the fresh
prime-$11$ variable and applying Theorem~\ref{thm:completion} proves
$n=11$.
This completes the proof of Theorem~\ref{thm:initial}.

\clearpage
\section{The middle intervals: 25 through 66}
\label{sec:finite}

\begin{theorem}[The middle intervals]\label{thm:finite}
For every field $K$, every $25\le n\le66$, and every choice
$c_a\in K^\times$, the polynomial
\[
 P=P_{n,c}=\sum_{a=1}^n c_aX(a)
\]
satisfies $d(P)=n$.
\end{theorem}

The constant multiplier gives $d(P)\le n$. For the lower bound we
first control small multipliers, then use residue descent to reduce
the remaining cases to them.

After rescaling the variables belonging to primes in $(n/2,n]$, write
\begin{equation}\label{eq:fin-setup}
 P=C+Z,\qquad Z=z_1+\cdots+z_r,\qquad
 r=\pi(n)-\pi(\lfloor n/2\rfloor),\quad c=\wt(C)=n-r,
 \quad R=\binom{r+1}{2}.
\end{equation}
The remaining variables are the base variables.

\subsection{Small multipliers}

First put
\begin{equation}\label{eq:fin-d0}
 d_0=1+\pi(\lfloor\sqrt n\rfloor)
       +\sum_{\substack{\sqrt n<q\le n/2\\q\text{ prime}}}
                       \left\lfloor\frac nq\right\rfloor.
\end{equation}
Then $d(C)\ge d_0$. Indeed, each prime variable $x_q$ with
$\sqrt n<q\le n/2$ occurs linearly, and two such variables never occur
together. Its coefficient is a weighted interval of length
$\lfloor n/q\rfloor\le6$, whose distance is that length by the initial
theorem. Repeated endpoint slicing adds these distances to the lower
bound $1+\pi(\lfloor\sqrt n\rfloor)$ for the remaining smooth core.
That lower bound follows from its constant and pure-prime vertices.

\begin{lemma}[Six input terms]\label{lem:fin-six}
For $25\le n\le66$, over every field and with arbitrary nonzero
weights, $\wt(PD)\ge n$ whenever $1\le\wt(D)\le6$.
\end{lemma}
\begin{proof}
A monomial input gives exactly $n$ terms.  Otherwise choose a
counterexample minimal under term deletion and normalize it.  If it has
$u$ terms and largest support prime $q$, the directed cut gives
$q<2(u-1)\le10$.  Thus $q\in\{2,3,5,7\}$, and the respective
minimum possible values of $u$ are $3,3,4,5$.

Let $C_q(t)$ count integers through $t$ with no prime factor at most
$q$. Split the output by its $q$-rough part, using the component
bounds of Lemma~\ref{arith:lem:active-prime-profile}.

For $q=2$, every component costs at least two, so the output has at
least $2C_2(n)=2\lceil n/2\rceil\ge n$ terms.  For $q=3$, the
lower bound is
\[
 3C_3(n)+C_3(n/2)-C_3(n/3).
\]
The residues $1,5$ modulo six give $3C_3(n)\ge n-1$.
Their successive gaps are at most four, while $(n/3,n/2]$ has length
greater than four.  The remaining difference is therefore at least one.

For $q=5$, every component costs at least four.  Write $n=30h+s$,
$0\le s<30$, and let $R(s)$ count the members of
$\{1,7,11,13,17,19,23,29\}$ not exceeding $s$.  The eight gaps give
$4R(s)-s\ge-2$.  Thus
\[
 4C_5(n)-n=2h+4R(s)-s\ge0\qquad(n\ge30).
\]
For $25\le n\le29$, the same residue list gives $4C_5(n)\ge n$.

For $q=7$, the successive component costs are
$5,4,6,6,8,8,5,5,\ldots$.  Their sum is
\[
 H(n)=5C_7(n)-C_7(n/2)+2C_7(n/3)+2C_7(n/5)-3C_7(n/7).
\]
Here are six short intervals on which monotonicity proves $H(n)\ge n$.
For $a\le n\le b$, evaluate $C_7(n/d)$ at $a/d$ in terms with
positive coefficients and at $b/d$ in terms with negative coefficients.
The resulting lower bound is $B$:
\begin{center}
\begin{tabular}{rr|rrrrr|r}
$a$&$b$&$C_7(a)$&$C_7(b/2)$&$C_7(a/3)$&$C_7(a/5)$&$C_7(b/7)$&$B$\\\hline
25&28&6&3&1&1&1&28\\
29&33&7&3&1&1&1&33\\
34&38&8&5&2&1&1&38\\
39&45&9&5&3&1&1&45\\
46&54&11&6&3&1&1&54\\
55&66&13&8&4&2&1&66
\end{tabular}
\end{center}
The $7$-rough integers through $66$ are
\[
 1,11,13,17,19,23,29,31,37,41,43,47,53,59,61,
\]
obtained by removing multiples of seven from the eight residues modulo
thirty. In every row $B=b$.
\end{proof}

The two minimum columns in the finite budget table
\eqref{eq:fin-block-budgets} give
\begin{equation}\label{eq:fin-basic-budgets}
 d_0+7r\ge n,\qquad rd_0>n.
\end{equation}
If $D$ uses only base variables, then
\[
 \wt(PD)=\wt(CD)+r\wt(D)\ge n.
\]
For at most six terms this is Lemma~\ref{lem:fin-six}; for at least
seven it follows from \eqref{eq:fin-basic-budgets}. We shall call this
the base-variable bound.

\subsection{Unique products and the odd-characteristic affine bound}

In this subsection assume $\operatorname{char}K\ne2$. Let $A$ be the
integer support of $C$, and define
\[
 \begin{gathered}
 S=\{ab:a,b\in A\},\\
 U=\{m\in S:\#\{(a,b)\in A^2:a\le b,\ ab=m\}=1\}.
 \end{gathered}
\]
Put $g=|U|$. At each position in $U$, the square coefficient is
$c_a^2$ or $2c_ac_b$, hence nonzero. Thus $\wt(C^2B)\ge g$ for
every monomial $B$.

\begin{lemma}[Lines parallel to a binomial]\label{lem:binomial-lines}
Let $F$ be a nonzero polynomial and let $B$ have exactly two terms.
If the exponent support of $F$ occupies $h$ affine lines parallel to
the difference of the two exponents of $B$, then $\wt(FB)\ge2h$.
\end{lemma}
\begin{proof}
On each occupied line, the two extreme terms of the product have
unique contributing pairs and nonzero coefficients. Different lines
remain disjoint after multiplication by either term of $B$.
\end{proof}

We use two elementary observations about the number of occupied lines.
A support of affine dimension $d$ occupies at least $d$ lines in any
fixed nonzero direction. If it contains four distinct collinear points,
and the chosen direction is not parallel to that line, it occupies at
least $d+2$ lines. Indeed, projection along the direction has affine
dimension at least $d-1$. The four points give four distinct collinear
images; at least $d-2$ further points are needed to obtain that dimension.

\begin{lemma}[Four small square bounds]\label{lem:finite-binomial-squares}
Assume $\operatorname{char}K\ne2$, and let $C$ be the weighted core
in \eqref{eq:fin-setup}. For every binomial $B$, including one involving
additional variables,
\[
 \wt(C^2B)\ge
 \begin{cases}
 16,&n=26,\\
 44,&n=27,28,\\
 20,&n=40.
 \end{cases}
\]
\end{lemma}
\begin{proof}
Write $x=x_2$. First consider $n=26$ or $40$. The support of $C^2$
has affine dimension $d=6$ or $8$, respectively: its constant term and
its extreme pure powers in the $d$ prime variables are nonzero.
The pure-$x$ part of $C$ is a dense polynomial of degree $e=4$ or $5$.
Its square has nonzero terms at exponents $0,1,2e-1,2e$. Thus, in
every direction other than the $x$-direction, it occupies at least
$d+2$ lines. In the $x$-direction the following odd integer positions
already give $d+2$ distinct occupied lines:
\[
 \begin{array}{c|l}
 n=26&1,3,5,7,11,13,49,121,\\
 n=40&1,3,5,7,11,13,17,19,49,121.
 \end{array}
\]
These square coefficients cannot vanish: the prime positions have
unique products $1\cdot p$, and $49=7^2$ and $121=11^2$ have unique
products in the core. Lemma~\ref{lem:binomial-lines} proves the bounds
$16$ and $20$.

Now let $n=27$ or $28$. Put $u=x_{11}$ and $v=x_{13}$. Then
\[
 C=A+uL+vM,
\]
where $A$ uses only $x_2,x_3,x_5,x_7$ and $L,M$ are nondegenerate
binomials in $x$. The six $(u,v)$-slices of $C^2$ are
\[
 A^2,\qquad 2AL,\qquad 2AM,\qquad L^2,\qquad 2LM,\qquad M^2.
\]
We show that their union occupies at least $22$ lines in every direction.

First, $A^2$ has at least $22$ terms. To see this,
set $x=0$, and write $y=x_3,z=x_5,w=x_7$. The resulting polynomial is
\[
 A_{\mathrm{odd}}=Q(y)+zR(y)+a z^2+wS(y),
\]
where $Q$ is a dense cubic, $R,S$ are nondegenerate binomials, and
$a\ne0$. In its square, the lower bounds for the $z$-row weights at
successive $w$-degrees are
\[
 \begin{array}{c|c|c}
 \text{$w$-degree}&\text{$z$-row lower bounds}&\text{total}\\\hline
 0&(4,2,1,2,1)&10\\
 1&(2,2,2)&6\\
 2&(3)&3.
 \end{array}
\]
Here $Q^2$ has nonzero endpoint and adjacent-endpoint coefficients;
$R^2+2aQ$ has a nonzero cubic coefficient; all other bounds follow
from endpoints. Hence $\wt(A_{\mathrm{odd}}^2)\ge19$. The additional
pure-$x$ square positions $2,128,256$ are nonzero, giving
$\wt(A^2)\ge22$.

If the direction has a nonzero component outside the four variables
of $A$, the terms of $A^2$ already lie on distinct parallel lines.
Otherwise the direction lies in those four variables, so the six
$(u,v)$-slices remain separate.

In the $x$-direction, $AL$ and $AM$ each occupy nine lines, indexed
by the nine odd parts
\[
 1,3,5,7,9,15,21,25,27.
\]
Each such row survives multiplication by the nonzero polynomial
$L$ or $M$. The other four slices each occupy at least one line,
so there are at least $9+9+4=22$ lines in total.

In a direction not parallel to $x$, the support of $A^2$ has affine
dimension four and contains four pure-$x$ points, so occupies at
least six lines. Both $AL$ and $AM$ have affine dimension four and
occupy at least four lines each. Finally, $L^2,M^2$ each have three
pure-$x$ terms, and $LM$ has at least two. These give distinct lines
in the chosen direction. The total is at least
\[
 6+4+4+3+2+3=22.
\]
Lemma~\ref{lem:binomial-lines} now gives $\wt(C^2B)\ge44$.
\end{proof}

\begin{lemma}[Two missing products]\label{lem:two-missing-products}
For $n=27,28$, every positive rational translate of the core support
$A$ covers at most $c-2$ unique-product positions:
\[
 |U\cap tA|\le c-2\qquad(t>0).
\]
\end{lemma}
\begin{proof}
Write $t=a/b$ in lowest terms. If $b\ge2$, an integer output requires
an input divisible by $b$, so there are at most
$\lfloor n/2\rfloor\le c-2$ such outputs. Assume henceforth that
$t$ is an integer.

If $t>n$, the inputs $n,n-1\in A$ both give outputs outside $A^2$.
Indeed, $tn>n^2$. For $t\ge n+2$ also $t(n-1)>n^2$; for $t=n+1$
this output is $n^2-1$, strictly between the largest two possible
products $n(n-1)$ and $n^2$.

If $1\le t\le n$ but $t\notin A$, it is a removed prime larger
than $n/2$ and divides no product in $A^2$. For $t\in A$ it remains
to exhibit two inputs giving nonunique products. At $t=1$ use
$4,6$; at $t=2$ use $2,3$. Both $4$ and $6$ have two factorizations.
For $3\le t\le n/2$, use the inputs $2,4$ and the distinct pairs
\[
 2t=1\cdot(2t)=2\cdot t,
 \qquad 4t=2\cdot(2t)=4\cdot t.
\]
Finally, if $n/2<t\le n$ and $t\in A$, then $t$ is composite.
Write $t=ps$ with $p$ its smallest prime factor. Here $p\le5$,
$s\ge4$, and $p,s,2p\in A$. The inputs $1,2$ give the distinct pairs
\[
 t=1\cdot t=p\cdot s,
 \qquad 2t=2\cdot t=(2p)\cdot s.
\]
In every case at least two inputs fail to give unique products.
\end{proof}

For a monomial $B$ and any polynomial $E$, the two-missing-products
lemma gives
\begin{equation}\label{eq:fin-unique-cover-bound}
 \wt(C^2B-CE)\ge g-(c-2)\wt(E)\qquad(n=27,28).
\end{equation}
Indeed each monomial translate of $C$ can cover at most $c-2$ of
the $g$ uniquely represented square positions. A translate involving
an incompatible extra variable covers none.

\begin{lemma}[All affine cofactors]\label{lem:fin-affine}
For $25\le n\le66$ and $\operatorname{char}K\ne2$, every nonzero
multiple $(C+Z)(D+ZB)$ with $D,B$ in the base polynomial ring has weight
at least $n$.
\end{lemma}
\begin{proof}
The exact weight is
\begin{equation}\label{eq:fin-affine-weight}
 m=\wt(CD)+r\wt(D+CB)+R\wt(B).
\end{equation}
The cases $B=0$ and $D=0$ follow from the base-variable bound; for the
latter use $r\ge1$ and $R\ge r$. Otherwise a strict counterexample
would give
\[
 t=\wt(B)\le\left\lfloor\frac{n-1-d_0}{R}\right\rfloor\le2,
 \qquad
 \wt(D+CB)\le K_t:=\left\lfloor\frac{n-1-d_0-Rt}{r}\right\rfloor.
\]
Writing $E=D+CB$, and using $CD=CE-C^2B$, gives
\begin{equation}\label{eq:fin-affine-budget}
 m\ge\wt(C^2B)-(n-2r)K_t+Rt.
\end{equation}
For the case $t=1$, it is enough to show
\[
 g\ge T_n:=n+(n-2r)K_1-R.
\]
Let $A_k$ be $[k]$ with the primes in $(k/2,k]$ removed, and put
\[
 f_k(m)=\#\{d\mid m:d\le\sqrt m,\ d\in A_k,\ m/d\in A_k\}.
\]
The sets $A_k$ increase with $k$: composites stay included, and a
prime already at most $k/2$ stays included. For $a\le n\le b$, a
product represented in $A_a^2$ and uniquely represented in $A_b^2$
remains uniquely represented in $A_n^2$. Consequently
\[
 g\ge\ell(a,b):=\sum_{m=1}^{a^2}
       \mathbf1_{\{f_a(m)>0,\ f_b(m)=1\}}.
\]
In the following table all minima and maxima are over $a\le n\le b$.
The prime counts, \eqref{eq:fin-d0}, and the displayed divisor formula
determine its entries by integer divisions.
\begingroup\small
\begin{equation}\label{eq:fin-block-budgets}
\begin{array}{c|r|r|r|r|r}
[a,b]&\min(d_0+7r-n)&\min(rd_0-n)&
 \max\left\lfloor\dfrac{n-1-d_0}{R}\right\rfloor&\max T_n&\ell(a,b)\\\hline
{[25,26]}&6&7&2&60&61\\
{[27,27]}&5&6&2&84&85\\
{[28,28]}&5&8&2&88&88\\
{[29,33]}&10&18&1&42&64\\
{[34,37]}&8&26&1&82&104\\
{[38,50]}&7&34&2&94&94\\
{[51,66]}&8&63&1&129&134
\end{array}
\end{equation}
\endgroup
The final two columns prove $g\ge T_n$. For a singleton interval,
$\ell(n,n)=g$; thus $g=85,88$ at $n=27,28$.

Only $n=26,27,28,40$ permit $t=2$. Apply
Lemma~\ref{lem:finite-binomial-squares} to
\eqref{eq:fin-affine-budget}:
\[
\begin{array}{r|rrrr}
n&\text{square--binomial bound}&K_2&R&\text{output bound}\\\hline
26&16&0&6&28\\
27&44&1&6&35\\
28&44&1&6&34\\
40&20&0&10&40
\end{array}
\]
Every final entry is at least $n$.
\end{proof}

In characteristic three the lemma remains valid when $D$ and $B$ also
use the fresh cubes $z_i^3$. Decompose them by those fresh-cube exponent
vectors. Output aggregate degrees zero, one, and two occupy disjoint
fresh residue classes, so the nonzero slices are ordinary affine
multiples and their weights add.

\subsection{Residue reduction and the higher-degree exceptions}

In characteristic three, an exponent class has integer representatives
$ab^3$ with a fixed cube-free part $a$. Hence
$\sigma_3(C)\le\lfloor n^{1/3}\rfloor\le r$ for $25\le n\le66$.
Similarly $\sigma_3(C^2)\le\lfloor n^{2/3}\rfloor\le R$ except
possibly at $n=26,27,28,38,39,40$, as follows by comparing
$n^2$ with $(R+1)^3$.

At $n=26$, a nonzero cube-residue class has $a\ge2$ and at most
$\lfloor(26^2/2)^{1/3}\rfloor=6$ terms. The zero class consists of
cubes $b^3$ with $b\le8$; neither $7^3$ nor $8^3$ is a product of
two integers at most $26$, so it also has at most six terms.
At $38\le n\le40$, nonzero classes have at most
$\lfloor(40^2/2)^{1/3}\rfloor=9$ terms. The zero class has $b\le11$
and excludes $7^3,11^3$, again giving at most nine, hence at most
$R=10$. The cubes $7^3,8^3,11^3$ have most balanced factor pairs
$(7,49),(16,32),(11,121)$, respectively.
Thus $\sigma_3(C^2)\le R$ except at $27,28$.
Corollary~\ref{cor:finite-residue} reduces every other
characteristic-three case to an affine cofactor, already handled
by Lemma~\ref{lem:fin-affine}.

For characteristic $p\ge5$, put $w_j=\binom{r+j-1}{j}$ and let $J$ be
the least positive index with $w_J\ge n$. The prime counts give $r=3$
only at $n=26,27,28$, and $r\ge4$ otherwise. For the two bounds
$(r,n)=(3,28)$ and $(4,66)$, the values of $w_1,\ldots,w_6$ are
\[
 (3,6,10,15,21,28),\qquad (4,10,20,35,56,84).
\]
Direct integer comparisons therefore give $J\le6$ and, for every
$1\le j<J$,
\begin{equation}\label{eq:fin-higher-residue}
 n^j\le w_j^{p_0},\qquad
 p_0=\begin{cases}5,&j\le4,\\7,&j=5.\end{cases}
\end{equation}
For $j<p$, a fixed exponent class modulo $p$ in the support of $C^j$
contains at most $\lfloor n^{j/p}\rfloor$ points: its integer
representatives have the form $a b^p\le n^j$. Thus
\eqref{eq:fin-higher-residue} gives $\sigma_p(C^j)\le w_j$ for $j<J$,
while $J\le j<p$ gives $\sigma_p(C^j)<n\le w_j$.
The residue theorem applies for every $p\ge5$.
In characteristic zero use scalar compression directly, with the same
increasing aggregate costs $w_j$.

The coefficient of the least aggregate row of any $P$-multiple is a
nonzero $C$-multiple: evaluate at $Z=-C$, divide by its least power of
$-C$, and reduce modulo $C$. By $rd_0>n$, a short output has a nonzero
constant row, costing at least $d_0$.
The inequality $d_0+w_j<n$ leaves only
$(n,j)=(26,3),(27,3),(27,4),(28,3),(28,4),(40,3)$.
After the constant and top rows, the respective remaining budgets
are $4,5,0,5,0,0$, giving the following possibilities:
\begin{center}
\begin{tabular}{c|c|>{\raggedright\arraybackslash}p{82mm}}
$n$&Top degree&Remaining possible rows\\\hline
$26,27,28$&3&Top row monomial, degree-two row absent, degree-one row
has at most one term.\\
$27,28$&4&Top row monomial, all intermediate rows absent.\\
$40$&3&Top row monomial, all intermediate rows absent.
\end{tabular}
\end{center}
These cases occur only in characteristic zero or $p\ge5$.

At $n=26,27,28$, set all base variables except those of primes
$5,7,11,13$ equal to zero. The remaining core is
$Q=A(x)+y_1+y_2+y_3$, after rescaling, where $A=a+bx+cx^2$ and
$abc\ne0$. In these characteristics $A^2$ has nonzero coefficients
in degrees $0,1,3,4$, while $A^3$ has nonzero coefficients in degrees
$0,1,5,6$. Therefore
\begin{equation}\label{eq:fin-cubic-bound}
 \wt(C^3)\ge\wt(Q^3)\ge4+3\cdot4+6\cdot3+10=44.
\end{equation}
For a cubic output the constant row is a translate of $C^3$ possibly
modified by one translate of $C$. Including the possible linear row
and the top fresh row gives the lower bounds
$44-c+3+10=34,33,32$, respectively. If the linear row is absent the
bound is stronger. In the quartic cases the degree-four part in
$y_1,y_2,y_3$ alone has fifteen terms; the top fresh row costs another
fifteen, giving $30>n$.
At $n=40$, retaining primes $7,11,13,17,19$ gives a weighted six-term
affine simplex. Its cube has $\binom83=56$ terms, and the top fresh
row adds twenty. These cases are also impossible.

\subsection{The two exceptional cube-residue cases}

It remains to handle $n=27,28$ in characteristic three. Here
$r=3$, $R=6$, $d_0=11,12$, and $\sigma_3(C^2)\le7$.
Indeed, nonzero classes have at most $\lfloor(28^2/2)^{1/3}\rfloor=7$
terms; the zero class has $b\le9$ and excludes $7^3,8^3$.
We handle these two cases directly.

Choose a short multiple of least total base cofactor width and compress. Theorem~\ref{thm:residue-descent} gives
\[
 G=A_0+ZA_1+Z^2A_2,\qquad A_j\in K[x,Z^3].
\]
With expanded weights $u_j=\wt(A_j)$, put
\[
 F=3u_1+6u_2,\qquad M=3u_1+7u_2=F+u_2.
\]
Then $m=\wt(G)=u_0+F<n$ and $u_0\ge d_0$. Hence $u_2\le2$;
if $u_2=1$ then $u_1\le3$, and if $u_2=2$ then $u_1\le1$.
The comparison and correction are
\begin{align*}
 V&=P(A_1-CA_2+ZA_2)
     =CA_1-C^2A_2+ZA_1+Z^2A_2,\\
 U&=G-V=A_0-CA_1+C^2A_2=P^3E.
\end{align*}
Every nonzero correction class of $U$ has at least $n$ terms by the
width descent. Its intersection with the comparison part is at most
$M$. There can be at most one such class, since
\[
 2(n-M)-(m-F)\ge n+1-F-2u_2\ge d_0+2-2u_2>0.
\]
If $u_2=0$, then $M=F$ and all classes are protected, so $U=0$.
If $U=0$ in any case, Lemma~\ref{lem:fin-affine}, applied to fresh-cube
slices, is already a contradiction. Otherwise $u_2=1$ or $2$ and
exactly one correction class is nonzero.

The comparison $V$ has at least $n$ terms. It also has at
least $2M$. If $u_2=1$, use \eqref{eq:fin-unique-cover-bound}:
\[
 \wt(V)\ge g-(c-5)u_1+6\ge34,
 \qquad 2M\le2(3\cdot3+7)=32.
\]
The two endpoint calculations are
$85-19\cdot3+6=34$ and $88-20\cdot3+6=34$.
If $u_2=2$, Lemma~\ref{lem:finite-binomial-squares} gives
$\wt(C^2A_2)\ge44$. If the two monomials have different
fresh-cube exponents, their contributions are disjoint and
the bound is stronger. Since $u_1\le1$,
\[
 \wt(V)\ge44-(c-3)u_1+12\ge34,
 \qquad 2M\le2(3+14)=34.
\]
Thus one correction class changes support by at least $n-2M$,
giving $m\ge\wt(V)+n-2M\ge n$. We have proved
\begin{proposition}\label{prop:fin-odd}
Theorem~\ref{thm:finite} holds in characteristic zero and every odd
characteristic.
\end{proposition}

\subsection{Characteristic two: stronger smooth-core bounds}

For the rest of this section, $\operatorname{char}K=2$.
A stronger smooth-core bound will sharpen $d_0$.
If $Q(t)=q_0+\cdots+q_et^e$ has $e>0$ and $q_0q_e\ne0$, endpoints give
$d(Q)\ge d(q_0)+d(q_e)$. If $Q$ cannot divide a two-row polynomial
over the fraction field of its coefficient ring, then every multiple
has a third occupied row and
\begin{equation}\label{eq:fin-third-row}
 d(Q)\ge d(q_0)+d(q_e)+1.
\end{equation}

A useful obstruction to a two-row multiple is
\begin{equation}\label{eq:fin-binomial-invariant}
 \frac{q_{e-1}^2}{q_eq_{e-2}}\text{ is nonconstant},
\end{equation}
when the three coefficients are nonzero. Indeed, in a splitting field
the roots of a binomial are $\rho$ times roots of unity, with their
multiplicities. A divisor therefore has
$q_{e-j}/q_e=\rho^j a_j$, where each $a_j$ is algebraic over $K$.
The displayed ratio would be algebraic over $K$. A rational function
algebraic over $K$ is constant: choose any variable on which it depends
and apply the one-variable argument in Lemma~\ref{lem:initial-quadratic}.
Thus a nonconstant ratio excludes a two-row multiple.

Let
\[
 T_{d,e}(x,y)=A_d(x)+yB_e(x)+c y^2,
 \qquad (d,e)\in\{(3,1),(3,2),(4,2)\},
\]
where $A_d,B_e$ are dense polynomials of the indicated degrees and
all displayed coefficients are nonzero. Each $T_{d,e}$ has distance
at least four. Its bottom row costs two, its top row one, and the
ratio $B_e^2/(cA_d)$ is nonconstant. For the first two pairs the degrees
$2e,d$ differ; for $(4,2)$ the numerator has no odd exponent while the
denominator does. Equation~\eqref{eq:fin-third-row} supplies the
fourth term.

Let $B_n$ denote the part of $C$ supported on integers with all prime
factors at most $\sqrt n$. We claim
\begin{equation}\label{eq:fin-binary-smooth}
 d(B_n)\ge b(n),\qquad
 b(n)=\begin{cases}
 7,&n=25,26,\\
 6,&27\le n\le48,\\
 8,&49\le n\le66.
 \end{cases}
\end{equation}
As before, each support below carries arbitrary nonzero coefficients.

For $n=25,26$, slice in $y=x_3$, with $x=x_2$ and $z=x_5$.
The bottom row is $T_{4,2}(x,z)$, of distance at least four; the top
row is a dense binomial in $x$, of distance two. The middle row has
support $1+x+x^2+x^3+z$. Its square has no odd $z$-exponent, whereas
the product of the bottom and top rows has a nonzero $z$-linear
coefficient. The quadratic ratio is nonconstant. Thus
$d(B_n)\ge4+2+1=7$.

For $27\le n\le66$, slice in $x_2$ and put
$y=x_3,z=x_5,w=x_7$. The last three row supports, in increasing
degree, are
\begin{center}
\begin{tabular}{c|c}
$n$&Last three $x_2$-row supports\\\hline
$27$--$31$, $48$--$55$&$(1+y+z,\ 1+y,\ 1)$\\
$32$--$39$, $64$--$66$&$(1+y,\ 1,\ 1)$\\
$40$--$47$&$(1+y+z,\ 1,\ 1)$\\
$56$--$63$&$(1+y+z+w,\ 1+y,\ 1)$
\end{tabular}
\end{center}
In each row of the table, the square of the middle coefficient has
only even exponents,
while the product of the other two coefficients has an odd exponent.
Thus~\eqref{eq:fin-binomial-invariant} forces a third occupied row.

For $27\le n\le48$, the bottom coefficient is $T_{3,1}(y,z)$
through $n=44$ and $T_{3,2}(y,z)$ thereafter, of distance at least four.
The top coefficient is monomial, so $d(B_n)\ge4+1+1=6$.

For $49\le n\le66$, the bottom row is
\[
 T=T_{3,2}(y,z)+wH(y,z)+cw^2,
\]
where $H$ has support $1+y+z$ through $n=62$ and $1+y+y^2+z$
thereafter. Since $H^2$ has only even exponents and $T_{3,2}$ has
nonzero odd exponents, the quadratic invariant is nonconstant.
Therefore $d(T)\ge4+1+1=6$.
The top $x_2$-row is monomial. The same odd-exponent obstruction in
the preceding table gives a third row, hence
$d(B_n)\ge6+1+1=8$. This proves \eqref{eq:fin-binary-smooth}.
Adding the prime rows above $\sqrt n$ by endpoints, exactly as for
$d_0$, now gives
\begin{equation}\label{eq:fin-kappa}
 d(C)\ge\kappa(n):=b(n)+
       \sum_{\substack{\sqrt n<q\le n/2\\q\text{ prime}}}
                         \left\lfloor\frac nq\right\rfloor.
\end{equation}

\subsection{Comparison bounds from four rough counts}

For a normalized $u$-term multiplier with largest support prime $q$,
summing the bounds of Lemma~\ref{arith:lem:active-prime-profile} gives:
\begin{align*}
 H_{2,u}(n)&=uC_2(n)+(2-u)C_2(n/2),\\
 H_{3,u}(n)&=uC_3(n)+(4-u)C_3(n/2)-C_3(n/3),\\
 H_{5,u}(n)&=uC_5(n)+(4-u)C_5(n/2)
                  +2C_5(n/3)-2C_5(n/5),\\
 H_{7,u}(n)&=uC_7(n)+(4-u)C_7(n/2)
                  +2C_7(n/3)+2C_7(n/5)-3C_7(n/7).
\end{align*}
These expressions follow from \eqref{arith:eq:threshold-sum}.
At $q=5$, for example,
the successive
costs are $u,4,6,4$ on indices $1$, $2$, $3$--$4$, and $5$ onwards.

For $27\le n\le48$, the component with rough part one at $q=5$
is exactly the smooth core $B_n$.  Its distance is at least six,
whereas $H_{5,u}$ charged four.  We may therefore add two to
$H_{5,u}$ throughout this range.  Denote the improved expression by
$H^+_{5,u}$, and use $H^+_{5,u}=H_{5,u}$ outside this range.

For $u\in\{3,4,5\}$, if a $u$-term multiplier uses a prime at least
$q_0$, the least prime at least $2u-1$, the directed cut gives
\begin{equation}\label{eq:fin-quantitative-cut}
 \wt(PD)\ge 2n-2(u-1)\lfloor n/q_0\rfloor.
\end{equation}
Indeed each proper piece has at most $u-1\le4$ terms and satisfies
Lemma~\ref{lem:fin-six}.  Otherwise its largest support prime is
one of $2,3,5,7$, and the corresponding bound above applies.

The resulting bounds are listed below. Each entry holds throughout its
row, and the last column is the minimum of the preceding entries.
\begin{center}\small
\begin{tabular}{c|r|r|r|r|r|r|r}
$n$&$u$&Cut&$H_{2,u}$&$H_{3,u}$&$H^+_{5,u}$&$H_{7,u}$&Bound\\\hline
25&3&30&33&28&---&---&28\\
27,28&4&32&42&33&32&---&32\\
34,35&4&40&50&40&40&---&40\\
36&4&42&54&44&40&---&40\\
38&4&46&56&48&44&---&44\\
39,40&4&48&60&47&46&---&46\\
49--51&4&56&74&62&60&---&56\\
58&5&76&100&78&73&69&69
\end{tabular}
\end{center}
The entries follow from
\[
 C_2(t)=\lfloor t\rfloor-\lfloor t/2\rfloor,
 \qquad C_3(t)=\lfloor t\rfloor-\lfloor t/2\rfloor
                    -\lfloor t/3\rfloor+\lfloor t/6\rfloor,
\]
the residue formula~\eqref{arith:residue-count}, and
\[
 C_7(t)=C_5(t)-C_5(t/7).
\]
For example, at $n=36$ the improved $q=5$ bound is
$4C_5(36)+2C_5(12)-2C_5(36/5)+2=36+6-4+2=40$.

\subsection{The only remaining four-term pattern}

\begin{lemma}[Parity pairing]\label{lem:fin-parity-pair}
In characteristic two, suppose a four-term Laurent multiplier has
three terms in one full exponent-parity class and its fourth term in
a different class.  Its weighted interval product has at least
$36,36,44$ terms at $n=27,28,36$, respectively.
\end{lemma}
\begin{proof}
For a weighted interval kernel $Q_k$ of length $k$ and a three-term
multiplier $F$, the initial theorem gives
\[
 \wt(Q_kF)\ge a(k),\qquad
 (a(0),a(1),\ldots,a(6))=(0,3,2,3,4,6,6).
\]
Only $a(5)=6$ needs explanation.  The strict case of sharp completion
applied to the dense quadratic core of $Q_5$ says that a five-term
multiple is a scalar monomial times $Q_5^q$, with $q$ a power of two.
Its cofactor is monomial if $q=1$; if $q\ge2$, its cofactor is a
nonzero $Q_5$-multiple and has at least five terms.  It cannot have
three terms.

After removing a Laurent monomial, write
$D=F^2+X(v)M^2$ over a coefficient field in which the needed square
roots exist.  Here $F$ has three terms, $M$ is monomial, and $v>1$
is squarefree.  For squarefree integers set
$s\oplus v=sv/\gcd(s,v)^2$.  The kernel splits into its parity classes:
\[
 P=\sum_{s\text{ squarefree}}X(s)Q_s^2,
 \qquad k_s:=\wt(Q_s)=\lfloor\sqrt{n/s}\rfloor,
\]
where $Q_s$ is a weighted interval of length $k_s$, and $k_s=0$
when $s>n$.  In parity class $s$, a square root of the output is
$Q_sF$ plus a monomial translate of $Q_{s\oplus v}M$.
For a pair with sizes $k,l>0$ the two classes consequently contribute
at least
\[
 \max\{a(k)-l,0\}+\max\{a(l)-k,0\}.
\]
If just one class is occupied, the contribution is $a(k)+k\ge4$.

If $v\ne2$, choose an odd prime divisor $p$ of $v$.  Every occupied
pair has exactly one member divisible by $p$.  The squarefree counts
through $27,28,36$ are $17,17,23$, respectively.  At most $4,4,6$
of those squarefree integers can be divisible by a fixed odd prime:
for $p=3$ the lists are $3,6,15,21$ and, at $36$, also $30,33$;
for $p=5$ the respective counts are $3,3,5$; for $p\ge7$ use
$\lfloor n/7\rfloor$.  Hence there are at least $9,9,11$ unpaired
occupied classes, giving $36,36,44$ output terms.

For $v=2$, pair each odd squarefree $s$ with $2s$:
\begin{center}
\begin{tabular}{r|l|r}
$n$&Pair-size multiset&Lower bound\\\hline
27&$(3,5),(2,3),(1,2),3(1,1),5(0,1)$&38\\
28&$(3,5),(2,3),2(1,2),2(1,1),5(0,1)$&36\\
36&$(4,6),(2,3),2(1,2),4(1,1),7(0,1)$&51
\end{tabular}
\end{center}
Here $j(k,l)$ means $j$ copies; each bound is sufficient.
\end{proof}

\subsection{Two largest parity classes suffice}

Suppose a short multiple exists. Choose one with least total base
cofactor width and apply residue descent:
\[
 G=A_0+ZD,\qquad U=A_0+CD=P^2E,\qquad V=PD,
 \qquad m=\wt(G)=\wt(A_0)+ru<n.
\]
Here $u=\wt(D)\ge1$, and every nonzero parity class of $U$ has at
least $n$ terms.

In $25\le n\le36$ one has $r\ge3,\kappa\ge12$; in
$37\le n\le66$ one has $r\ge4,\kappa\ge18$.  Thus
$r\kappa\ge n$.  The residue theorem gives
\begin{equation}\label{eq:fin-binary-u}
 u\le\left\lfloor\frac{n-1-\kappa}{r}\right\rfloor.
\end{equation}
Using~\eqref{eq:fin-kappa} gives the following table. Its last column
lists those $n$ for which both
$\kappa+4r<n$ and $2(3\lfloor\sqrt n\rfloor+
\lfloor\sqrt{n/2}\rfloor)>n$ hold.
\begin{center}
\begin{tabular}{c|r|l}
Range&$\min(\kappa+5r-n)$&Possible four-term exceptions\\\hline
25--26&3&none\\
27--33&1&27,28\\
34--40&1&34,35,36,38,39,40\\
41--48&4&none\\
49--57&1&49,50,51\\
58&$-1$&none\\
59--66&0&none
\end{tabular}
\end{center}
Consequently $u\le4$ except possibly at $n=58$, where
$(r,\kappa)=(6,27)$ gives $u\le5$.  Lemma~\ref{lem:fin-six}
therefore supplies $\wt(V)\ge n$ in every case.

Put
\[
 s=\lfloor\sqrt n\rfloor,\qquad t=\lfloor\sqrt{n/2}\rfloor.
\]
These are exactly the two largest parity-class sizes of $C$: a
squarefree class $a$ contains the integers $ab^2\le n$, hence has
$\lfloor\sqrt{n/a}\rfloor$ elements; removing fresh primes affects
only singleton classes.  If the parity distribution of $D$ is
$\lambda_1\ge\lambda_2\ge\cdots$, the comparison $CD$ has at most
$M=\sum_i\lambda_i s_i$ terms in any one class, where $s_i$ are
the decreasing class sizes of $C$.  In particular,
\begin{equation}\label{eq:fin-majorization}
 M\le us;\qquad
 M\le(u-1)s+t\quad\text{if $D$ meets at least two classes}.
\end{equation}
For $u=4$ meeting at least two classes, unless the distribution is
$(3,1)$, the stronger bound
$M\le2s+2t$ holds.  These follow simply by putting as many terms
as allowed against the largest class and all remaining terms
against the second largest.

If $a$ correction classes are nonzero, residue descent gives
\begin{equation}\label{eq:fin-capacity-gates}
 m\ge\wt(V)+a(n-2M),\qquad
 a(n-M)\le m-ru.
\end{equation}
Thus $2M\le n$ immediately excludes a counterexample.  If
$2M\le n+ru$, the second inequality forces $a\le1$; in that
case the first excludes a counterexample whenever $\wt(V)\ge2M$.

Single-class inputs with $u\ge3$ can be excluded directly.  Let $Q_1$ count squarefree integers in $(n/4,n]$.
Their parity classes in $P$ are singletons, so each contributes
exactly $u$ terms to $V$.  Fresh-prime classes cannot meet $U$.
Every other such class is protected: a correction would leave at
least $n-u>m-ru$ terms in $A_0$.  Hence $m\ge uQ_1$.
For $n\le66$, inclusion--exclusion for the squares $4,9,25,49$
gives the exact formula
\[
 Q_1=n-2\lfloor n/4\rfloor-\lfloor n/9\rfloor
       -\lfloor n/25\rfloor-\lfloor n/49\rfloor
       +\lfloor n/16\rfloor+2\lfloor n/36\rfloor.
\]
Since $1/25+1/49<1/16$, this is greater than $4n/9-3$.
Thus $3Q_1\ge n$ for $27\le n\le66$; at $n=25,26$, one has
$Q_1=11,12$, respectively.  Single-class inputs with $u\ge3$ are impossible.

For $u\le2$, the first bound in~\eqref{eq:fin-majorization} gives
$2M\le4s\le n$.  For $u=3$ meeting two classes,
$2M\le4s+2t\le n$ for every $n\ge26$ in the range.  Indeed the
real-variable inequality holds for $n\ge30$, and $26\le n\le29$
has $(s,t)=(5,3)$.  At $n=25$, $M\le13$ and the comparison bound
$28\ge2M$ applies, while $2M\le n+3r$.

For $u=4$, the preceding count table leaves only its stated
four-term exceptions.  All have $2(3s+t)\le n+4r$, so there is
at most one correction class.  Their capacities and comparison
bounds reduce to
\begin{center}
\begin{tabular}{c|r|r|r|r}
$n$&$3s+t$&$2s+2t$&General bound&$(3,1)$ bound\\\hline
27,28&18&16&32&36\\
34,35&19&18&40&40\\
36&22&20&40&44\\
38&22&20&44&44\\
39,40&22&20&46&46\\
49--51&$\le26$&$\le24$&56&56
\end{tabular}
\end{center}
The only improvements in the last column are the three cases of
Lemma~\ref{lem:fin-parity-pair}.  Each bound is at least twice the
corresponding capacity, completing the four-term case.

Finally, if $u=5$, then $n=58$.  A mixed parity distribution has
$M\le4\cdot7+5=33$, and $2M\le58+5\cdot6$.  The comparison bound
$\wt(V)\ge69>66\ge2M$ excludes this last possibility.

\begin{proof}[Proof of Theorem~\ref{thm:finite}]
The preceding argument proves characteristic two.
Proposition~\ref{prop:fin-odd} proves characteristic zero and odd characteristic.
The monomial multiplier attains the stated bound in every case.
\end{proof}

\section{Arithmetic estimates}\label{sec:arithmetic}

Throughout, the interval coefficients are nonzero elements of a field.
Write $r(n)=\pi(n)-\pi(n/2)$.
For a finite set $T$ of primes, let $C_T(x)$ count the positive
integers through $x$ divisible by no prime in $T$, and put
$J_T(x)=C_T(x)-C_T(x/2)$. Write $C_q,J_q$ when $T$ consists of
all primes at most $q$. An integer is \emph{$q$-smooth} if all its
prime factors are at most $q$, and \emph{$q$-rough} if none are.

\subsection{Disjoint components}

Let $D$ be normalized, with $u$ terms and largest support prime $q$.
Multiplication by $D$ preserves the $q$-rough factor. Hence
\begin{equation}\label{arith:eq:prime-rank-decomposition}
 \wt(P_{n,c}D)=\sum_{\substack{b\le n\\b\text{ is }q\text{-rough}}}
                    \wt(S^{(b)}_{q,\lfloor n/b\rfloor}D),
\end{equation}
where $S^{(b)}_{q,j}$ has the $q$-smooth integers through $j$ as
support and inherits the coefficients $c_{ba}$. These coefficients
may vary with $b$.

Component costs can be summed by thresholds. For any cost function
$f$ on $\{1,\ldots,n\}$,
\begin{equation}\label{arith:eq:threshold-sum}
 \sum_{\substack{b\le n\\b\text{ is }q\text{-rough}}}
 f(\lfloor n/b\rfloor)
 =f(1)C_q(n)+\sum_{j=2}^n\bigl(f(j)-f(j-1)\bigr)C_q(n/j).
\end{equation}
Indeed, $f(k)=f(1)+\sum_{j=2}^k(f(j)-f(j-1))$, and a rough part
has component index at least $j$ exactly when $b\le n/j$.
This identity explains the coefficients in the counting bounds below.

\begin{lemma}[Component bounds]\label{arith:lem:active-prime-profile}
The component of index $j=1$ has weight $u$, and for $j\ge2$,
\[
 \wt(S^{(b)}_{q,j}D)\ge
 \begin{cases}
 2(1+\pi(j)),&2\le j<q,\\
 1+\pi(q),&j\ge q.
 \end{cases}
\]
Consequently, for $q\ge5$ and $u\ge4$,
\begin{equation}\label{arith:eq:four-component-budget}
 \wt(P_{n,c}D)\ge4C_q(n)+(u-4)J_q(n),
\end{equation}
and for $q\ge17$,
\begin{equation}\label{arith:eq:eight-component-budget}
 \wt(P_{n,c}D)\ge uJ_q(n)+4C_q(n/2)+2C_q(n/3)+2C_q(n/5).
\end{equation}
\end{lemma}
\begin{proof}
At index one the kernel is a scalar. If $2\le j<q$, it is
independent of $x_q$, so at least two nonzero $x_q$-slices of $D$
remain separate. Each costs at least $1+\pi(j)$ terms by the
dimension bound. If $j\ge q$, the kernel contains the constant and
a pure power of every prime through $q$, giving $1+\pi(q)$ terms.
Applying \eqref{arith:eq:threshold-sum} to the lower profiles
$u,4,4,\ldots$ and $u,4,6,6,8,8,\ldots$ gives the two inequalities.
\end{proof}

\begin{theorem}[Localized multipliers]\label{arith:thm:localized-expansion}
Let $n\ge67$, and let $D$ be a normalized nonconstant multiplier
with $u$ terms and largest support prime $q$. If
\[
 q\le n/2,\qquad q<2(u-1),
\]
then $\wt(P_{n,c}D)\ge n$ over every field.
\end{theorem}

The proof is completed in Section~\ref{arith:sec:localized-completion}.

\subsection{Sieve estimates}\label{arith:sec:one-step-sieve}

The eight residues
$\mathcal R=\{1,7,11,13,17,19,23,29\}$ give the exact formula
\begin{equation}\label{arith:residue-count}
 C_5(x)=\sum_{a\in\mathcal R}\left\lfloor\frac{x+30-a}{30}\right\rfloor.
\end{equation}
Equivalently, if $\lfloor x\rfloor=30h+t$ with $0\le t<30$,
then $C_5(x)=8h+\#\{a\in\mathcal R:a\le t\}$.

\begin{lemma}[Sieve discrepancies]\label{arith:lem:rough-discrepancy}
\label{arith:lem:no-core-half-discrepancy}
Put $\rho(q)=\prod_{p\le q}(1-1/p)$. For every $x\ge0$,
\[
 |C_q(x)-\rho(q)x|\le\Delta_q,\qquad
 |J_q(x)-\rho(q)x/2|\le2\Delta_q,
\]
where
\[
 \Delta_2=\tfrac12,\qquad \Delta_3=\tfrac23,\qquad
 \Delta_q=\tfrac{14}{15}2^{\pi(q)-3}\quad(q\ge5).
\]
For $q=5$ the sharper half-interval bound is
\[
 |J_5(x)-2x/15|\le22/15.
\]
\end{lemma}
\begin{proof}
The cases $q=2,3$ follow from the residues modulo two and six.
At the eight jumps of $C_5$ in one period, the values of
$15C_5(x)-4x$ are
\[
 11,2,1,8,7,14,13,4.
\]
Between jumps the error $C_5(x)-4x/15$ decreases linearly; its
left limit at a jump is one less than its value. Thus its absolute
value is at most $14/15$.
For a prime $p\notin T$,
\[
 C_{T\cup\{p\}}(x)=C_T(x)-C_T(x/p),
\]
which multiplies the density by $1-1/p$ and at most doubles the
error. This proves the bound for $C_q$; subtraction at $x$ and
$x/2$ gives the bound for $J_q$.

The error $J_5(x)-2x/15$ has period sixty. Between jumps it has
slope $-2/15$, and at its upward jumps its numerator is
\[
\begin{array}{c|rrrrrrrr}
t&1&7&11&13&17&19&23&29\\\hline
15J_5(t)-2t&13&1&8&19&11&22&14&2\\[2pt]
t&31&37&41&43&47&49&53&59\\\hline
15J_5(t)-2t&13&1&-7&4&-4&7&14&2
\end{array}
\]
Thus the error is at most $22/15$. For nonintegral $0<x<60$,
reflection of the residues gives $J_5(60-x)=8-J_5(x)$; hence the
error changes sign. One-sided limits give the lower bound at
jumps as well.
\end{proof}

\begin{lemma}[A one-step sieve]\label{arith:lem:one-step-sieve}
Let $T$ and $P$ be disjoint finite sets of primes, and suppose
$|C_T(x)-\rho x|\le\Delta$ for all $x\ge0$. Put
\[
 \alpha=\rho\left(1-\sum_{p\in P}\frac1p\right),
 \qquad E=(1+|P|)\Delta.
\]
Then
\[
 C_{T\cup P}(x)\ge\alpha x-E,\qquad
 J_{T\cup P}(x)\ge\alpha x/2-2E.
\]
If the sharper estimate $|J_T(x)-\rho x/2|\le\Lambda$ is known,
the second error can be replaced by $(1+|P|)\Lambda$.
\end{lemma}
\begin{proof}
Removing multiples of the added primes gives
\[
 C_{T\cup P}(x)\ge C_T(x)-\sum_{p\in P}C_T(x/p).
\]
Multiplication by $p\notin T$ preserves the $T$-rough property.
Applying the same bound on $(x/2,x]$ gives
\[
 J_{T\cup P}(x)\ge J_T(x)-\sum_{p\in P}J_T(x/p).
\]
Insert the assumed discrepancies, using $\Lambda=2\Delta$ if no
sharper bound is available.
\end{proof}

For $u\ge4$ and $q\ge5$, combining this lemma with
\eqref{arith:eq:four-component-budget} gives
\begin{equation}\label{arith:eq:one-step-component-bound}
             \wt(P_{n,c}D)\ge\frac{u+4}{2}\alpha n-2(u-2)E,
\end{equation}
where $P$ consists of the primes at most $q$ outside the chosen
sieve $T$.

For $T=\{2,3,5\}$ the constants are
\begin{equation}\label{arith:eq:base-five-sieve}
 \alpha_q=\frac4{15}\left(1-\sum_{7\le p\le q}\frac1p\right),
 \qquad E_q=\frac{14}{15}\bigl(\pi(q)-2\bigr).
\end{equation}
For $T=\{2,3,5,7\}$, use $\rho=8/35$ and $\Delta=28/15$.

\subsection{Finite prime counts}\label{arith:sec:written-low-range}
\label{sec:prime-blocks}

\begin{lemma}[Prime counting below a rough cube]
\label{arith:lem:rough-semiprime-count}
Put
\[
 \begin{gathered}
 C_{11}(x)=C_5(x)-C_5(x/7)-C_5(x/11)+C_5(x/77),\\
 (p_1,\ldots,p_8)=(13,17,19,23,29,31,37,41).
 \end{gathered}
\]
For $11\le x\le1799$,
\begin{equation}\label{arith:eq:small-prime-sieve}
 \pi(x)=4+C_{11}(x)
       -\sum_{i=1}^8\max\{C_{11}(x/p_i)-i,0\}.
\end{equation}
\end{lemma}
\begin{proof}
Inclusion--exclusion gives the formula for $C_{11}$. Since
$1799<13^3$, an eleven-rough integer in this range is one, a prime,
or a product $pq$ of two primes with $13\le p\le q$.
The smaller prime is one of the listed $p_i$, since
$\sqrt{1799}<43$. Also $1799/13<13^2$, so $C_{11}(x/p_i)$ counts
only one and primes. Exactly $i$ eleven-rough integers lie below $p_i$.
Thus the summand counts the composites with smaller prime $p_i$,
including squares. Subtract these and one, then add the five
primes through eleven.
\end{proof}

\begin{lemma}[A prime budget]\label{arith:lem:eight-endpoint-prime-budget}
Let $D$ be a normalized multiplier at $n$, of weight
$u\ge U\ge4$, where $U$ is an integer, and largest support prime
$17\le q\le n/2$. Suppose $q<2(u-1)$ and
$r=\pi(n)-\pi(n/2)\ge8$. Then
\[
 \wt(P_{n,c}D)\ge Ur+4\pi(n/2)+2\pi(n/3)+2\pi(n/5)
                         +8-8\pi(2U-3).
\]
At $U=10$ this becomes
\begin{equation}\label{arith:eq:ten-prime-budget}
 \wt(P_{n,c}D)\ge10r+4\pi(n/2)+2\pi(n/3)+2\pi(n/5)-48.
\end{equation}
\end{lemma}
\begin{proof}
Every prime in $(n/2,n]$ is $q$-rough, so $J_q(n)\ge r$.
Also $C_q(t)\ge1+\pi(t)-\pi(q)$ for all $t\ge0$: when the right
side is positive it counts $1$ and primes larger than $q$; otherwise
the assertion is immediate. The component budget
\eqref{arith:eq:eight-component-budget} therefore gives
\[
 \wt(P_{n,c}D)\ge ur+4\pi(n/2)+2\pi(n/3)+2\pi(n/5)-8\pi(q)+8.
\]
The hypothesis gives $q\le2u-3$. Each unit increase of the integer
parameter in $2u-3$ adds at most one new odd prime. Thus
\[
 \pi(q)\le\pi(2U-3)+(u-U).
\]
Substitution leaves the nonnegative term $(u-U)(r-8)$.
For $U=10$, use $\pi(17)=7$.
\end{proof}

\begin{lemma}[The counts through 1023]
\label{arith:lem:first-prime-budget-block}\label{arith:prop:written-low-range}
For $67\le n\le1023$,
\[
 r(n)\ge\lfloor\sqrt n\rfloor,\qquad
 10r(n)+4\pi(n/2)+2\pi(n/3)+2\pi(n/5)-48\ge n.
\]
\end{lemma}
\begin{proof}
All counts below follow from \eqref{arith:eq:small-prime-sieve}.
First suppose $67\le n\le100$. At $n=67$ the primes in $(n/2,n]$ are
$37,41,43,47,53,59,61,67$. Through $n=100$, primes enter at
$71,73,79,83,89,97$ and leave at $74,82,86,94$.
Pairing each departure with the preceding arrival gives
\[
 r(n)\ge8+\mathbf1_{n\ge73}+\mathbf1_{n\ge97}
       \ge\lfloor\sqrt n\rfloor.
\]
Also $\pi(n/2)\ge11$, $\pi(n/3)\ge8$, and $\pi(n/5)\ge6$,
so the second expression is at least $104\ge n$.

For $101\le n\le1023$, a row $[a,b]$ gives the lower bounds
\[
 r_0=\pi(a)-\pi(b/2),\qquad
 B_0=10r_0+4\pi(a/2)+2\pi(a/3)+2\pi(a/5)-48.
\]
Monotonicity gives $r(n)\ge r_0$ and a prime budget at least $B_0$.
Each row satisfies $r_0\ge\lfloor\sqrt b\rfloor$ and $B_0\ge b$:
\[
\begin{array}{rr|rr}
a&b&r_0&B_0\\\hline
101&117&10&150\\
118&141&11&172\\
142&165&12&200\\
166&195&13&228\\
196&225&15&262\\
226&273&16&302\\
274&345&19&354\\
346&421&22&430\\
422&524&27&524\\
525&636&33&636\\
637&765&40&772\\
766&913&48&916\\
914&1023&59&1102
\end{array}
\]
For example, the row $[422,524]$ uses prime counts
$82,55,47,34,23$ at $422,262,211,140,84$, giving
$r_0=27$ and $B_0=524$.
\end{proof}

For $67\le n\le228$, let $\kappa,h$ be as in
Section~\ref{sec:large-weighted}. On each block $[a,b]$ below put
\[
 r_0=\pi(a)-\pi(b/2),\quad R_0=r_0(r_0+1)/2,\quad
 \kappa_0=\kappa(a),\quad h_0=\pi(a/2)-\pi(\sqrt b).
\]
Then $r\ge r_0$ and $h\ge h_0$. Moreover $\kappa\ge\kappa_0$:
between consecutive prime squares the set of primes at most $\sqrt n$
is fixed, the existing floor terms increase, and each new prime entering
at $n=2q$ adds two. The only prime squares in this range are $121,169$,
both block starts. The finite sieve formula and integer division give
the following table; both final columns are at least $b$.
\[
\begin{array}{rr|rrr|rr}
a&b&r_0&\kappa_0&h_0&\kappa_0+R_0+r_0&\kappa_0+R_0+h_0+1\\\hline
67&72&8&28&7&72&72\\
73&81&9&30&7&84&83\\
82&85&9&37&9&91&92\\
86&93&9&40&10&94&96\\
94&101&9&47&11&101&104\\
102&115&10&50&11&115&117\\
116&120&13&58&12&162&162\\
121&133&12&53&12&143&144\\
134&145&12&62&14&152&155\\
146&161&12&71&16&161&166\\
162&168&14&80&17&199&203\\
169&193&15&72&17&207&210\\
194&223&15&88&19&223&228\\
224&228&18&109&23&298&304
\end{array}
\]
This proves \eqref{eq:large-endpoint-budgets}.

Put $d_0=1+\pi(n/2)$ and $R=r(r+1)/2$. For the next table use
$r_0=\pi(a)-\pi(b/2)$, $d_0=1+\pi(a/2)$ and
$R_0=r_0(r_0+1)/2$.
The same endpoint bounds give $d_0+R\ge n$ for $229\le n\le1799$:
\[
\begin{array}{rr|rrrr}
a&b&r_0&d_0&R_0&d_0+R_0\\\hline
229&241&20&31&210&241\\
242&262&21&31&231&262\\
263&286&22&33&253&286\\
287&325&24&35&300&335\\
326&361&25&39&325&364\\
362&397&27&43&378&421\\
398&457&29&47&435&482\\
458&541&31&51&496&547\\
542&654&34&59&595&654\\
655&837&39&67&780&847\\
838&1093&45&82&1035&1117\\
1094&1480&52&102&1378&1480\\
1481&1799&80&132&3240&3372
\end{array}
\]

\subsection{Fresh primes}\label{arith:sec:elementary-primes}

\begin{lemma}[Elementary Chebyshev bounds]
\label{arith:lem:elementary-chebyshev}
Put
\[
 A=\frac{14\log2+9\log3+5\log5}{30}>\frac{23}{25},
 \qquad b=\frac65A,
\]
\[
 H(x)=\frac5{4\log6}
       \left(\log x+\frac12\log6-\frac65\right)^2+9.
\]
For $x\ge30$,
\[
 Ax-\frac52\log x+2<\psi(x)< bx+H(x),
\]
where
\[
 \psi(x)=\sum_{p^a\le x}\log p,\quad \vartheta(x)=\sum_{p\le x}\log p.
\]
For $x\ge256$, moreover,
\[
                  0\le\psi(x)-\vartheta(x)<\frac83\sqrt x.
\]
Consequently, with $r(n)=\pi(n)-\pi(n/2)$,
\[
 r(n)\log n>
 \frac{2A}{5}n-\frac83\sqrt n-\frac52\log n+2-H(n/2)
 \qquad(n\ge256).
\]
\end{lemma}
\begin{proof}
Write $J(t)=\log\lfloor t\rfloor!$ and $h(t)=t\log t-t$.
Concavity of the logarithm, applied on the intervals between consecutive
integers and on the intervals centred at consecutive integers, gives
\[
 h(t)-\frac12\log t+\frac{1+\log2}{2}
 \le J(t)\le h(t)+\frac12\log t+1\qquad(t\ge1).
\]
The upper bound at an integer $k$ follows from the
trapezoidal lower bound for $\int_1^k\log t\,dt$; its extension to real
$t\ge1$ is monotonic. Jensen's inequality gives
$J(t)\ge\int_{1/2}^{\lfloor t\rfloor+1/2}\log v\,dv$.
Since $t<\lfloor t\rfloor+1$ and $h$ has derivative $\log t$,
\[
 h(t)-h(\lfloor t\rfloor+1/2)\le\tfrac12\log t.
\]
This proves the lower bound.

Consider the factorial combination
\[
 Q(x)=J(x)-J(x/2)-J(x/3)-J(x/5)+J(x/30).
\]
The displayed bounds imply, for $x\ge30$,
\[
 Ax-\frac52\log x+\log60-2
 \le Q(x)\le
 Ax+\frac52\log x-\log30+\frac{1-3\log2}{2}.
\]
In particular,
\begin{equation}\label{arith:eq:elementary-factorial-combination}
 Ax-\frac52\log x+2<Q(x)<Ax+\frac52\log x-3.
\end{equation}

The function
\[
 f(t)=\lfloor t\rfloor-\lfloor t/2\rfloor-\lfloor t/3\rfloor
                      -\lfloor t/5\rfloor+\lfloor t/30\rfloor
\]
is periodic with period $30$, and is constant on each unit interval.
In one period its value is one precisely on the intervals whose left
endpoints are
\[
             1,2,3,4,5,7,8,9,11,13,14,17,19,23,29;
\]
elsewhere its value is zero. In particular it is one on $[1,6)$.
The factorial valuation formula therefore gives
\[
 Q(x)=\sum_{p^a\le x}f(x/p^a)\log p,\qquad
                \psi(x)-\psi(x/6)\le Q(x)\le\psi(x).
\]
The lower bound for $\psi$ follows at once.

We also need the elementary bound $\psi(t)<t\log4$ for $t\ge1$.
Equivalently, $\operatorname{lcm}(1,\ldots,k)<4^k$ for integers $k\ge1$.
The case $k=1$ is immediate. Induct on $k$. For even $k$, the lcm increases from its preceding value
by a factor at most two. For $k=2m+1$, each new prime power above $m+1$
contributes its prime as a factor of $\binom{2m+1}{m}$, so
\[
 \operatorname{lcm}(1,\ldots,2m+1)
 \le\operatorname{lcm}(1,\ldots,m+1)\binom{2m+1}{m}
 <4^{m+1}4^m.
\]
The binomial bound follows by adding its two equal central coefficients.

Iterate $\psi(x)\le Q(x)+\psi(x/6)$ until
$t=x/6^k$ lies in $[5,30)$. The linear terms sum to $b(x-t)$, whereas
$\psi(t)<t\log4$ and $(\log4-b)t<9$. The remaining sum is at most
\[
 \sum_{j=0}^{k-1}\left(\frac52(\log x-j\log6)-3\right)
 =-\frac54\log6\,k^2+
       \left(\frac52\log x+\frac54\log6-3\right)k
 \le H(x)-9,
\]
by completing the square. This proves the upper bound for $\psi$.

To remove prime powers, absorb all even exponents into $\psi(\sqrt x)$:
\[
 \psi(x)-\vartheta(x)
 =\psi(\sqrt x)+\sum_{\substack{a\ge3\\a\text{ odd}}}
                       \vartheta(x^{1/a}).
\]
There are at most $(\log x)/(2\log2)-3/2$ possible odd exponents at
least five. The preceding elementary bound for $\psi$ thus gives
\[
 \frac{\psi(x)-\vartheta(x)}{\sqrt x}
 <\log4\left(1+x^{-1/6}+
       \left(\frac{\log x}{2\log2}-\frac32\right)x^{-3/10}\right).
\]
For $x\ge256$, the second term in parentheses is less than $2/5$.
The last term decreases on this range, since $\log2>2/3$, and at $256$
it equals $(5/2)2^{-12/5}<1/2$. Hence the parentheses are less than
$19/10$, and $\log4<7/5$ gives the claimed bound $8/3$.

Finally $r(n)\log n\ge\vartheta(n)-\vartheta(n/2)$; insert the two
estimates to obtain the stated bound. The logarithm comparisons used
here and below follow from
\[
 \log\frac{1+t}{1-t}=2\sum_{j\ge0}\frac{t^{2j+1}}{2j+1}
 \qquad(0<t<1),
\]
obtained by integrating the geometric series. The first three terms at
$t=1/3$ give $\log2>693/1000$, and bounding the remaining denominators
below by three gives $\log2<25/36<7/10$.
The first four terms at $t=1/2$ give $\log3>1098/1000$.
Finally $\log5=2\log2+\log(5/4)$; the first two terms at $t=1/9$
give $\log5>1609/1000$. These three lower bounds give $A>23/25$
and $\log6>25/14$. The same upper-tail bound at $t=1/5$ gives
$\log(3/2)<73/180<21/50$, so
$\log2+6/5-\tfrac12\log6>99/100$.
\end{proof}

\begin{lemma}[Fresh prime counts]\label{arith:lem:no-core-fresh-budget}
For every $n\ge67$, $r(n)\ge\lfloor\sqrt n\rfloor$. Moreover,
\[
 \begin{array}{ll}
 r(n)>n/30,&1024\le n\le32768,\\[2pt]
 r(n)>7n/(20\log n)>6\sqrt n,&n\ge32768.
 \end{array}
\]
\end{lemma}
\begin{proof}
Lemma~\ref{arith:prop:written-low-range} covers $67\le n\le1023$.
Assume $n\ge1024$, and put $L=\log n$ and
\[
 \begin{aligned}
 P(L)&=\frac7{10}\left(L-\frac{99}{100}\right)^2+\frac52L+7,\\
 K(L)&=\frac{46}{125}-\frac83e^{-L/2}-e^{-L}P(L).
 \end{aligned}
\]
Lemma~\ref{arith:lem:elementary-chebyshev} gives
\begin{equation}\label{arith:eq:no-core-K}
                         \frac{r(n)L}{n}>K(L).
\end{equation}
Indeed, its error term satisfies
$H(n/2)<(7/10)(L-99/100)^2+9$: use
$5/(4\log6)<7/10$,
$\log2+6/5-(1/2)\log6>99/100$, and $L>6$.

For $L\ge6$, $K$ is increasing and strictly concave.  Direct
differentiation reduces these assertions to $P-P'>0$ and
\[
 P-2P'+P''
 =\frac7{10}\left\{\left(L-\frac{299}{100}\right)^2-2\right\}
       +\frac52L+2>0.
\]
Both inequalities hold throughout $L\ge6$; for the first, write
$P-P'=(7/10)\{(L-99/100)^2-2(L-99/100)\}+(5/2)L+9/2$.
Thus $K(L)-L/30$ is strictly concave as well.

At $n=1024$, use $L<7$, $\sqrt n=32$, and $P(7)<50$ to obtain
\[
 K(L)-L/30>
 \frac{46}{125}-\frac1{12}-\frac{50}{1024}-\frac7{30}
 =\frac{481}{192000}>0.
\]
At $n=32768=2^{15}$, use $L<21/2$, $\sqrt n>180$, and
$P(21/2)<97$ to obtain
\[
 K(L)-\frac7{20}>
 \frac{46}{125}-\frac2{135}-\frac{97}{32768}-\frac7{20}
 =\frac{24881}{110592000}>0.
\]
Since $L/30<7/20$ at this endpoint, concavity proves $K(L)>L/30$
throughout the first range.  Since $K$ is increasing, its value remains
greater than $7/20$ throughout the second range.  Equation
\eqref{arith:eq:no-core-K} proves both estimates.

For $1024\le n\le32768$, $n/30>\sqrt n$. For larger $n$,
$\sqrt n/\log n$ is increasing and, at $32768$,
\[
 \frac7{20}\frac{\sqrt n}{\log n}
       >\frac7{20}\frac{180}{21/2}=6.
\]
\end{proof}

\subsection{Small support primes}
\label{arith:sec:small-prime-profile-replacement}

\begin{proposition}[Small support primes]
\label{arith:prop:small-prime-profile-replacement}
Let $n\ge67$. If a normalized nonconstant $D$ has $u$ terms and largest support prime
$q\le13$ with $q<2(u-1)$, then $\wt(P_{n,c}D)\ge n$.
\end{proposition}

\begin{proof}
For $q=2$, every rough component has at least two terms, so
\[
 \wt(P_{n,c}D)\ge2C_2(n)=2\lceil n/2\rceil\ge n.
\]
For $q=3$, the hypothesis gives $u\ge3$.  Each component
has at least three terms except that the index-two components have
at least four: their two nonempty $3$-valuation slices each have
two endpoints.  Thus
\[
 \wt(P_{n,c}D)\ge3C_3(n)+C_3(n/2)-C_3(n/3).
\]
The residues $1,5$ modulo $6$ give $3C_3(n)\ge n-1$.
Their successive gaps are at most four, so the interval
$(n/3,n/2]$, of length $n/6\ge4$, contains a $3$-rough integer.
This proves the required bound for $q=3$.

For $q=5$, every component has at least four terms, giving
\[
 \wt(P_{n,c}D)\ge4C_5(n)
 \ge\frac{16n}{15}-\frac{56}{15}\ge n
 \qquad(n\ge56).
\]
For $q=7,11,13$, the hypothesis gives $u\ge5,7,8$, respectively.
The component profile gives
\begin{align*}
 H_7(n)&:=5C_7(n)-C_7(n/2)+C_7(n/3),\\
 H_{11}(n)&:=7C_{11}(n)-3C_{11}(n/2)+2C_{11}(n/3),\\
 H_{13}(n)&:=8C_{13}(n)-4C_{13}(n/2)
             +2C_{13}(n/3)+C_{13}(n/5),
\end{align*}
and $\wt(P_{n,c}D)\ge H_q(n)$.  Indeed, after the index-one component
of weight $u$, the lower profiles are respectively
\[
 (4,5,5,\ldots),\qquad (4,6,6,\ldots),\qquad
 (4,6,6,7,7,\ldots),
\]
starting at component index two. Increasing $u$ above these minima
adds the corresponding nonnegative multiple of $J_q(n)$.

Equivalently,
\begin{align*}
 H_7&=4C_7+J_7+C_7(n/3),\\
 H_{11}&=4C_{11}+3J_{11}+2C_{11}(n/3),\\
 H_{13}&=4C_{13}+4J_{13}+2C_{13}(n/3)+C_{13}(n/5).
\end{align*}
The base-$5$ one-step sieve has constants
\[
\begin{array}{c|ccc}
q&7&11&13\\\hline
\alpha_q&8/35&236/1155&184/1001\\
E_q&28/15&14/5&56/15
\end{array}
\]
where $C_q(t)\ge\alpha_qt-E_q$ and
$J_q(t)\ge\alpha_qt/2-2E_q$.  Substituting in the positive
expressions above gives
\[
 H_7(n)\ge\frac{116n-1372}{105},\qquad
 H_{11}(n)\ge\frac{4366n-116424}{3465},\qquad
 H_{13}(n)\ge\frac{18952n-840840}{15015}.
\]
These bounds exceed $n-1$ for $n\ge116,126,210$, respectively,
so integrality gives the result in those ranges. The remaining
ranges are covered by the following blocks.  For a row $[a,b]$ and a column $q$, put
\[
 \beta_7=\min(b,115),\quad
 \beta_{11}=\min(b,125),\quad \beta_{13}=b.
\]
When $a\le\beta_q$, monotonicity of $C_q$ gives the lower bound
$B_q$ obtained from $H_q$ by evaluating terms with positive
coefficients at $a/d$ and the sole term with a negative coefficient
at $\beta_q/2$.  Thus $H_q(n)\ge B_q$ on the remaining part
$[a,\beta_q]$ of that block.
\[
\begin{array}{c|r|r|r}
[a,b]&B_7&B_{11}&B_{13}\\\hline
{[67,76]}&76&89&91\\
{[77,85]}&86&102&105\\
{[86,94]}&94&110&114\\
{[95,101]}&101&121&127\\
{[102,110]}&110&132&139\\
{[111,145]}&126&149&149\\
{[146,181]}&-&-&181\\
{[182,209]}&-&-&238
\end{array}
\]
Every entry is at least the upper endpoint $\beta_q$ to which
it applies.  The dashes mean that the preceding linear estimate
already applies.

The entries follow from \eqref{arith:residue-count} and
\[
 C_7(t)=C_5(t)-C_5(t/7),\quad
 C_{11}(t)=C_7(t)-C_7(t/11),\quad
 C_{13}(t)=C_{11}(t)-C_{11}(t/13).
\]
For example, in the row $[111,145]$ the counts at
$111,\beta_q/2,37,111/5$ are respectively
$(26,13,9,5)$, $(25,14,8,4)$, and $(24,15,7,3)$.
They give $B_7=126$, $B_{11}=149$, and $B_{13}=149$.
\end{proof}

\subsection{Larger support primes}\label{arith:sec:no-core-closure}

\begin{lemma}[Reciprocal sums]
\label{arith:lem:no-core-reciprocals}
One has
\[
 \sum_{7\le p\le53}\frac1p<\frac{259}{400},\qquad
 \sum_{11\le p\le149}\frac1p<\frac{229}{320},\qquad
 \sum_{11\le p\le271}\frac1p<\frac{33}{40}.
\]
\end{lemma}
\begin{proof}
Since $\sqrt{271}<17$, the primes from $17$ through $271$ are
exactly the integers in that range coprime to $2\cdot3\cdot5\cdot7\cdot11\cdot13$.
Adding $11,13$ gives the following list.
Let $B_j$ be the reciprocal sum of these primes in
$[30j,30j+30)$, truncated at $271$.  The row sums satisfy:
\[
\begin{array}{c|l|r}
j&\text{denominators in }B_j&4000B_j\text{ is less than}\\\hline
0&11,13,17,19,23,29&1429\\
1&31,37,41,43,47,53,59&657\\
2&61,67,71,73,79,83,89&381\\
3&97,101,103,107,109,113&230\\
4&127,131,137,139,149&147\\
5&151,157,163,167,173,179&146\\
6&181,191,193,197,199&105\\
7&211,223,227,229,233,239&106\\
8&241,251,257,263,269&79\\
9&271&15
\end{array}
\]
The first five upper bounds sum to $2844$, and all ten sum to $3295$.
The first budget follows from the first two rows, since
\[
 \sum_{7\le p\le53}\frac1p
 <\frac17+\frac{1429+657}{4000}-\frac1{59}
 =\frac{259}{400}-\frac{19}{206500}<\frac{259}{400}.
\]
The other two reciprocal sums are less than $711/1000$ and
$659/800$, respectively.  Since
$711/1000<229/320$ and $659/800<33/40$, the claims follow.
\end{proof}

For a prime $q\ge17$, put
\[
 h_q=\pi(q)-2,\qquad
 \alpha_q=\frac4{15}\left(1-\sum_{7\le p\le q}\frac1p\right).
\]
The one-step sieve, applied separately to $C$ and $J$ with their two
discrepancies, gives
\[
 C_q(x)\ge\alpha_qx-\frac{14}{15}h_q,
 \qquad
 J_q(x)\ge\frac{\alpha_qx}{2}-\frac{22}{15}h_q.
\]

\begin{lemma}[Support primes from 17 through 53]
\label{arith:lem:no-core-small-primes}
Let $D$ be a normalized multiplier of weight $u$, with
largest support prime $17\le q\le53$ and $q<2(u-1)$.
For every $n\ge1024$, one has $\wt(P_{n,c}D)>n$.
\end{lemma}
\begin{proof}
Rewrite \eqref{arith:eq:eight-component-budget} as
\[
 \wt(P_{n,c}D)\ge
 4C_q(n)+(u-4)J_q(n)+2C_q(n/3)+2C_q(n/5).
\]
The preceding discrepancies imply
\begin{equation}\label{arith:eq:no-core-positive-profile}
 \wt(P_{n,c}D)\ge
 \alpha_q\left(\frac u2+\frac{46}{15}\right)n
                      -\frac{22u+24}{15}h_q.
\end{equation}
Use the hypothesis $q<2(u-1)$ and the five columns
\[
\begin{array}{c|rrrrr}
q\text{ in}&[17,23]&[29,37]&[41,43]&\{47\}&\{53\}\\\hline
u_*&10&16&22&25&28\\
a&7/50&7/60&13/125&99/1000&47/500\\
h_*&7&10&12&13&14
\end{array}
\]
Here $u\ge u_*$, $h_q\le h_*$, and $\alpha_q>a$ throughout each
band.  The reciprocal comparisons all follow from the single inequality
\[
                  \sum_{7\le p\le53}\frac1p<\frac{259}{400},
\]
from Lemma~\ref{arith:lem:no-core-reciprocals}: subtract the terms
above the endpoint of each band.
At $n=1024$ and $u=u_*$, the right side of
\eqref{arith:eq:no-core-positive-profile}, with $a,h_*$ in place of $\alpha_q,h_q$,
exceeds $n$ by, respectively,
\[
 \frac{6964}{375},\quad\frac{10672}{225},\quad\frac{126832}{1875},\quad
 \frac{106162}{1875},\quad \frac{40192}{1875}.
\]
The slope in $u$ is at least
$1024(47/1000)-22\cdot14/15>27$.
The positive corner values also imply a positive slope in $n$.
Both slopes increase with the other variable, proving the bound
for all $n\ge1024$ and $u\ge u_*$. 
\end{proof}

\begin{lemma}[Support primes from 59 through 271]
\label{arith:lem:no-core-middle-primes}
Let $D$ be a normalized multiplier of weight $u$,
with largest support prime $59\le q\le271$ and $q<2(u-1)$.
For every $n\ge32768$, one has $\wt(P_{n,c}D)>n$.
\end{lemma}
\begin{proof}
Use $R=\{2,3,5,7\}$ as the fixed small sieve.  Its density and
discrepancy are $\rho=8/35$ and $\Delta=28/15$, the latter following
by adding $7$ to the modulo-thirty sieve.  The one-step bound is
\[
 \wt(P_{n,c}D)\ge\frac{u+4}{2}\alpha n-2(u-2)E,
\]
where
\[
 \alpha=\frac8{35}\left(1-\sum_{11\le p\le q}\frac1p\right),
 \qquad E=\frac{28}{15}\{\pi(q)-3\}.
\]
The modulo-thirty discrepancy gives
$\pi(q)\le2+C_5(q)<3+4q/15$, hence $E<112q/225$.

Lemma~\ref{arith:lem:no-core-reciprocals} gives the following bounds.
The last column is the resulting lower bound minus $n$ at
$n=32768$ and $u=u_*$:
\[
\begin{array}{c|rrrr}
q&u_*&\alpha>&E<&\text{margin}\\\hline
59\text{--}149&31&13/200&75&778/5\\
151\text{--}271&77&1/25&135&1654/25
\end{array}
\]
The slope in $u$ is at least $n/50-270>0$. The positive corner
values imply a positive slope in $n$, and both slopes increase
with the other variable. This proves the claim.
\end{proof}

\begin{lemma}[An elementary prime-product estimate]
\label{arith:lem:no-core-odd-product}
For every odd prime $q$,
\[
                \rho(q)=\prod_{p\le q}(1-1/p)>\frac1{2\sqrt q}.
\]
\end{lemma}
\begin{proof}
Including all odd composite factors decreases the product.  For every
odd integer $m\ge3$, one has $(1-1/m)^2>(m-2)/m$.  Multiplication
over $m=3,5,\ldots,q$ telescopes to give
\[
 \rho(q)\ge\frac12\prod_{\substack{3\le m\le q\\m\text{ odd}}}
                (1-1/m)>\frac1{2\sqrt q}.
\]
\end{proof}

\begin{lemma}[Support primes at least 277]
\label{arith:lem:no-core-large-primes}
Let $D$ be a normalized multiplier of weight $u$.
Suppose $n\ge32768$, $u<(20/7)\log n$, and its largest support
prime satisfies $277\le q<2(u-1)$. Then $\wt(P_{n,c}D)>5n/4$.
\end{lemma}
\begin{proof}
Put $L=\log n$. The hypotheses give
\[
 q<\frac{40}{7}L,\qquad
 L>\frac7{40}(q+2)\ge\frac{1953}{40}>48.
\]
The eleven-rough count has density $16/77$ and discrepancy $56/15$.
For $q\ge277$, it includes its unit, all primes greater than $11$ and at most $q$,
and the two composite integers $169=13^2$ and $221=13\cdot17$.
Consequently
\[
 \pi(q)\le C_{11}(q)+2
 \le\frac{16}{77}q+\frac{86}{15}
 <\frac{16}{77}q+6.
\]
For the full $q$-rough sieve, write
$\Delta_q=(14/15)2^{\pi(q)-3}$.  Since $\log2<7/10$,
\[
 \Delta_q<\frac{112}{15}
       \exp\left(\frac{16}{77}q\log2\right)
 <\frac{112}{15}\exp\left(\frac{64}{77}L\right).
\]
The finite error therefore satisfies
\[
 \frac{2u\Delta_q}{n}
 <\frac{128}{3}L e^{-13L/77}
 <2048e^{-624/77}<\frac{2048}{2800}<\frac34.
\]
Here $Le^{-13L/77}$ decreases for $L\ge48$, and
$624/77>8$ while $e^8=e^6e^2>20^2\cdot7=2800$.
Indeed,
\[
 e^3>\sum_{j=0}^{8}\frac{3^j}{j!}
       =\frac{89641}{4480}>20,
 \qquad
 e^2>\sum_{j=0}^{4}\frac{2^j}{j!}=7.
\]

On the other hand, odd-integer telescoping and $u>q/2$ give
\[
                      \frac{u\rho(q)}2>\frac{\sqrt q}{8}>2.
\]
The positive component budget yields $\wt(P_{n,c}D)\ge uJ_q(n)$, and
the unconditional rough discrepancy gives
$J_q(n)\ge\rho(q)n/2-2\Delta_q$.  Thus
\[
                \frac{\wt(P_{n,c}D)}n>2-\frac34=\frac54>1,
\]
as claimed.
\end{proof}

\subsection{Completion of the localized estimate}
\label{arith:sec:localized-completion}

\begin{proof}[Proof of Theorem~\ref{arith:thm:localized-expansion}]
Proposition~\ref{arith:prop:small-prime-profile-replacement} handles
$q\le13$. Otherwise $q\ge17$ and $u\ge10$. For $67\le n\le1023$,
Lemma~\ref{arith:prop:written-low-range} and
\eqref{arith:eq:ten-prime-budget} give the result.

Every prime in $(n/2,n]$ is $q$-rough and contributes an index-one
component, so $\wt(P_{n,c}D)\ge ur$. We may assume $u<n/r$.
For $1024\le n\le32768$, Lemma~\ref{arith:lem:no-core-fresh-budget}
gives $u<30$, hence $q\le53$, and
Lemma~\ref{arith:lem:no-core-small-primes} applies.
For $n\ge32768$, the same lemma gives $u<(20/7)\log n$.
Lemmas~\ref{arith:lem:no-core-small-primes},
\ref{arith:lem:no-core-middle-primes}, and
\ref{arith:lem:no-core-large-primes} cover respectively
$q\le53$, $59\le q\le271$, and $q\ge277$, which include every
remaining prime.
\end{proof}

\section*{Acknowledgements}
This work would not exist without the paper of Philippa Holdridge and
P\'eter P\'al Pach~\cite{HP}, whose framework it builds on throughout,
as described in Section~\ref{sec:main}. We are equally grateful to
Pilz, who posed the problem and settled its first cases, and to Huang,
Ke, Meyer, Nagy, Pach and Szab\'o, whose work shaped
it~\cite{Pilz,HKP,PS,KM,NP}. Any errors are ours. The formal verification
depends on Lean~4 and on mathlib, and we thank their developers and
contributors.

\paragraph{How this work was produced.}
The arguments in this paper, its written proofs, and the Lean~4
formalisation were developed with substantial assistance from AI
systems. The written proofs have not been checked line by line by a
human mathematician. Correctness rests on the formalisation: the main
theorems and corollaries, and the intermediate results recorded in the
proof map, are machine-checked using only the standard axioms of Lean's
foundations, and every module of the development has been re-checked by
an independent kernel replay (\texttt{lean4checker}). Where the printed
argument and the formal proof differ, the formal proof is authoritative.
We state this so that readers can weigh the work appropriately.

\clearpage
\begin{thebibliography}{9}
\raggedright
\bibitem{Pilz}
G. F. Pilz,
\emph{On polynomial near-ring codes},
in \emph{Contributions to General Algebra}, vol.~8 (Linz, 1991),
H\"older-Pichler-Tempsky, Vienna, 1992, pp.~233--238.
\bibitem{HKP}
P.-Y. Huang, W.-F. Ke and G. F. Pilz,
\emph{The cardinality of some symmetric differences},
Proc. Amer. Math. Soc. \textbf{138}(3) (2010), 787--797.
\href{https://doi.org/10.1090/S0002-9939-09-10189-2}
{doi:10.1090/S0002-9939-09-10189-2}.
\bibitem{PS}
P. P. Pach and C. Szab\'o,
\emph{On the minimal distance of a polynomial code},
Discrete Math. Theor. Comput. Sci. \textbf{13}(4) (2011), 33--44.
\href{https://doi.org/10.46298/dmtcs.556}
{doi:10.46298/dmtcs.556}.
\bibitem{KM}
W.-F. Ke and J. H. Meyer,
\emph{On a conjecture regarding the symmetric difference of certain sets},
Bull. Aust. Math. Soc. \textbf{111}(3) (2025), 397--404.
\href{https://doi.org/10.1017/S0004972724000807}
{doi:10.1017/S0004972724000807}.
\bibitem{HP}
P. Holdridge and P. P. Pach,
\emph{On the Extended 1-2-3 Conjecture of Pilz},
2026, \href{https://arxiv.org/abs/2607.00934}{arXiv:2607.00934}.
\bibitem{NP}
J. Nagy and P. P. Pach,
\emph{A sumset version of a conjecture of Pilz},
SIAM J. Discrete Math. \textbf{40}(2) (2026), 675--679.
\href{https://doi.org/10.1137/24M1654609}
{doi:10.1137/24M1654609};
\href{https://arxiv.org/abs/2409.15075}{arXiv:2409.15075}.
\bibitem{Archive}
\emph{Pilz's Extended 1--2--3 Conjecture: proofs and Lean~4
formalisations}, Zenodo, 2026. Contains this paper and its Lean
development, an earlier proof of the binary case, and a note on the
density bound. Concept DOI (resolving to the latest version):
\href{https://doi.org/10.5281/zenodo.23568997}{doi:10.5281/zenodo.23568997}.
\end{thebibliography}
\end{document}
